% Clase del documento:
\documentclass[a4paper,12pt]{article}


% Paquete geometry para establecer los márgenes:
\usepackage[inner=3.5cm, outer=3cm, top=4cm, bottom = 4cm,
            headheight=2cm, headsep=1.5cm, footskip=1.5cm]{geometry}

\usepackage[spanish,es-tabla]{babel}
\usepackage[utf8]{inputenc}
\usepackage[T1]{fontenc}


% Bibliografía

%MIGUEL A. %
%\usepackage{natbib}
%MIGUEL A. %
\usepackage[nottoc]{tocbibind}
\usepackage{doi}
%MIGUEL A. %\bibliographystyle{apacite}
%JAVI
\bibliographystyle{splncs04}


% Paquete para imprimir codigo:
\usepackage{listings}
\lstset{breaklines=true,
        language=C,
        extendedchars=true,
        basicstyle=\small}
\renewcommand{\lstlistingname}{Código}

% Paquetes para imprimir pseudocodigo:
\usepackage{algorithm}
\usepackage{algpseudocode}
\usepackage{algorithmicx}

% Traducción
\floatname{algorithm}{Algoritmo}

% Paquetes para ecuaciones y numeros:
\usepackage{amsmath}
\usepackage{amsthm}
\usepackage{amsfonts}
\usepackage{amssymb}
\newtheorem{defn}{Definición}
\newtheorem{thm}{Teorema}
\newtheorem{lem}{Lema}
\newenvironment{case1}{\textit{Caso 1:}}{}
\newenvironment{case2}{\textit{Caso 2:}}{}
\newenvironment{case3}{\textit{Caso 3:}}{}

% Paquete para espaciado e interlineados:
\usepackage{setspace}

% Paquete para incluir figuras y gráficos:
\usepackage{graphicx}
\usepackage{titlesec}
\usepackage{appendix}
\usepackage{caption}
\usepackage{subcaption}
\usepackage{makecell}
\usepackage{svg}
\usepackage{pdfpages}
\captionsetup{justification=raggedright, labelfont=bf}

% Estilo de enlaces y urls:
%\usepackage[usenames]{color}
\usepackage{hyperref}
\usepackage{url}

\hypersetup{
    colorlinks=true,
    linkcolor=black,
    urlcolor=blue,
%    citecolor=blue
}

% Estilo de paginas:
\usepackage{lmodern}
\usepackage{fancyhdr}
\usepackage{etoolbox}

\renewcommand\qedsymbol{$\blacksquare$}

\DeclareMathOperator*{\argmax}{arg\,max}
\DeclareMathOperator*{\argmin}{arg\,min}

%%%%%%%%%%%%%%%%%%%%%%%%%%%%%%%%%%%%%%%%%%%%%%%%%%%%%%%%%%%%%%%%%%%%%%%%%%%%%%%

\title{{\Large{Trabajo de Fin de Grado}} \\
        \vspace{1.25cm}
        {{\textbf{Estrategias de paralelización del algoritmo A* en entornos de memoria compartida}}}  \\





% {\color{blue} 
% Pon un título más concreto



% Como esquema genra, podrias ser algo mas o menos así:

% 1.- Introducción


% 2.- Estado del arte

% Donde intentes englobar todo lo que ya hay hecho en este campo, con referencias bibliograficas.

% Si ves que lo tienes separado en diferentes secciones, tambien
% puedes poner en lugar de un Estado de arte general, diferentes subsecciones con estados del arte de cada seccion de la que hables



% 3.- Análisis de objetivos y metodología
	
%     3.1 Objetivos
	
%     3.2 Metodología

% donde dices lo que pretendes (objetivos)
% y como lo quieres conseguir (metodologia)

% Un par de paginas en total


% 4 Diseño y solución

% Bibliografia

% }
        
        \vspace{1cm}
        \begin{figure}[htbp]
            \centering
	        \includegraphics[scale=0.25]{./images/logo-um.png}
        \end{figure}
}
\author{\Large{Universidad de Murcia} \\
    \\
    \\
    Autor:
    \\
    \large{Miguel Ángel Rodríguez García} \\
    \\
    \\
    Directores:
    \\
    \large{A. Javier Cuenca Muñoz y José Matías Cutillas Lozano}
    \\
    \\
}
\date{\Large{2025}}
%%%%%%%%%%%%%%%%%%%%%%%%%%%%%%%%%%%%%%%%%%%%%%%%%%%%%%%%%%%%%%%%%%%%%%%%%%%%%%%

% Comienzo del documento:
\begin{document}

% Se establece el interlineado del documento:
%\onehalfspacing

% Se establece la separación entre párrafos:
\setlength{\parskip}{0.15cm}

% Se crea la portada:
\maketitle
\thispagestyle{empty}


\newpage

\null\cleardoublepage

\newpage

%%%%%%%%%%%%%%%%%%%%%%%%%%%%%%%%%%%%%%%%%%%%%%%%%%%%%%%%%%%%%%%%%%%%%%%%%%%%%%%
\section*{Resumen}
%%%%%%%%%%%%%%%%%%%%%%%%%%%%%%%%%%%%%%%%%%%%%%%%%%%%%%%%%%%%%%%%%%%%%%%%%%%%%%%

En este trabajo exploramos las distintas estrategias de paralelización del algoritmo A* en entornos de memoria compartida. El tema de este trabajo ha estado motivado por la realización de un Máster de Computación Cuántica en la Universidad Internacional de La Rioja durante el actual curso académico. En dicho Máster se propuso una línea de Trabajo de Fin de Máster (TFM) enfocada en una propuesta cuántica del algoritmo A*. El planteamiento principal en este trabajo ha sido la paralelización de dicho algoritmo de cara a obtener una base sólida sobre la que comparar una propuesta cuántica en un futuro.

El algoritmo A* es un algoritmo de búsqueda heurístico utilizado para resolver el problema de encontrar el camino de menor coste entre dos nodos en un grafo. Este algoritmo ha probado ser una técnica muy eficiente que a día de hoy se aplica en una gran variedad de campos como el desarrollo de videojuegos, el transporte inteligente, la planificación robótica o la inteligencia artificial. Sin embargo, a la hora de buscar una formulación paralela surgen varias complicaciones, algunas inherentes a los entornos paralelos, y otras debidas al fuerte carácter secuencial propio del propio algoritmo.

Comenzaremos el trabajo con una revisión del tratamiento de este problema a lo largo de la historia. Veremos cuándo surgen sus primeras formulaciones y estrategias para resolverlo, y lo estudiaremos no solo desde un punto de vista de interés para resolver problemas actuales a gran escala, sino también como un problema para el que la naturaleza ha desarrollado mecanismos de adaptación eficientes para resolverlo. Presentaremos los fundamentos matemáticos necesarios para comprender tanto la formulación del problema del camino más corto como el funcionamiento del propio algoritmo A*. Entre éstos se incluirán conceptos de teoría de grafos como nodos, aristas, grafos dirigidos y ponderados, así como la noción de camino, coste asociado y camino óptimo. Expondremos formalmente la formulación del problema del camino más corto como la búsqueda de una secuencia de nodos que minimice el coste total desde un nodo de inicio hasta un nodo objetivo, bajo la condición de que cada nodo consecutivo esté conectado mediante una arista válida.

A continuación, revisaremos brevemente los distintos paradigmas de algoritmos de búsqueda que existen a día de hoy, clasificados en dos grandes grupos: informados y no informados. Los algoritmos no informados, como la búsqueda en profundidad (DFS), búsqueda en anchura (BFS) o búsqueda de coste uniforme (UCS), no disponen de información sobre la distancia estimada al objetivo, lo que les lleva a expandir un mayor número de nodos y, por tanto, a ser menos eficientes. En contraposición, los algoritmos informados o heurísticos, como el algoritmo primero el mejor (GBFS) o A*, utilizan una función heurística para guiar la búsqueda hacia la solución de forma más eficiente.

El algoritmo A* fue propuesto en 1968 por Hart, Nilsson y Raphael en el contexto del proyecto Shakey, uno de los primeros robots con capacidad de razonamiento sobre sus acciones. La principal aportación de A* fue la introducción formal del uso de la información heurística en los problemas de búsqueda. Para ello, introducen el uso de la esta información para estimar el coste restante hasta la meta y la combinan con el coste acumulado desde el inicio. La función de evaluación que proponen para A* y que guía la búsqueda es $f(n) = g(n) + h(n)$, donde $g(n)$ representa el coste desde el inicio hasta el nodo $n$ y $h(n)$ es el coste para alcanzar el objetivo desde $n$. Este acercamiento permite encontrar caminos óptimos de forma más eficiente que los algoritmos clásicos. Veremos cómo más allá de ser un único algoritmo, A* engloba en realidad a una familia de algoritmos, siendo la búsqueda de coste uniforme y la búsqueda primero el mejor casos particulares de A*. Mostraremos un ejemplo detallado de su funcionamiento para ilustrar el proceso de expansión de nodos y el uso de las listas abierta y cerrada.

También haremos mención a las diferentes adaptaciones que han surgido del algoritmo para adaptarlo tanto a dominios específicos como a necesidades particulares en cuanto a uso de memoria o entornos que pueden variar en tiempo de ejecución. Gran parte de estas variantes fueron motivadas por el principal problema que presenta A*, que es su gran consumo de memoria durante la ejecución.

Analizaremos también las propiedades fundamentales del algoritmo, como son la admisibilidad y la optimalidad. Veremos cómo estas propiedades están relacionadas con la naturaleza de la función heurística utilizada, dejando ver que la elección y el diseño de ésta es una de las partes más importantes del algoritmo. La propiedad de admisibilidad nos garantizará que A* encontrará la solución óptima siempre que la función heurística no sobreestime el coste real (es decir, sea optimista), mientras que la propiedad de optimalidad se refiere a que A* expande el menor número de nodos necesario para encontrar dicha solución óptima. En particular, mostraremos que si la heurística es consistente (también llamada monótona), el algoritmo no necesitará explorar varias veces el mismo nodo. Para el estudio de la optimalidad, haremos un análisis del artículo de Dechter y Pearl, en el que proponen distintos tipos de optimalidad para distintos tipos de algoritmos en función de las características del problema al que se enfrentan, y finalizaremos con una clasificación en la que se especifica bajo qué condiciones y para qué tipos de problemas A* es un algoritmo óptimo.

Una vez comprendido el funcionamiento secuencial de A*, nos centramos en el estado del arte de su paralelización. Durante las últimas décadas, se han propuesto múltiples estrategias para adaptar este algoritmo a entornos paralelos, especialmente motivadas por la necesidad de reducir los tiempos de ejecución en problemas de gran escala. Estas estrategias se pueden clasificar principalmente en dos enfoques: paralelización centralizada y descentralizada.

En las estrategias centralizadas, todos los procesos involucrados en la búsqueda comparten una lista abierta global y una lista cerrada común, que son las estructuras de datos principales que utiliza el algoritmo. La versión más representativa de este enfoque es SPA* (\textit{Simple Parallel A*}), donde los procesos compiten por extraer nodos de la lista abierta y expandirlos, sincronizándose mediante mecanismos de exlusión mutua para evitar condiciones de carrera. Esta estrategia es sencilla de implementar en entornos de memoria compartida, pero presenta importantes limitaciones en cuanto a escalabilidad. A medida que aumenta el número de procesos, la contención por acceder a las listas compartidas incrementa, y el tiempo de espera entre procesos puede superar la ganancia obtenida por la paralelización. Además, el criterio de terminación del algoritmo debe redefinirse para asegurar que se ha encontrado la mejor solución global posible, ya que varios procesos pueden estar expandiendo nodos en paralelo sin conocer el estado completo de la búsqueda. Veremos cómo este tipo de estrategias se han aplicado a dominios en los que la operación de expansión es costosa y la contención en las estructuras de datos compartidas representa una pequeña parte del tiempo de ejecución total del algoritmo.

Por otro lado, las estrategias descentralizadas reparten el trabajo entre procesos de forma que cada uno mantenga su propia lista abierta y cerrada. Esta estrategia introduce nuevos retos, como la detección de duplicados y la necesidad de sincronizar el estado de la búsqueda entre procesos. Sin embargo, permite una mayor escalabilidad y mejora el rendimiento en problemas de gran tamaño. Veremos distintas estrategias de reparto de trabajo que han surgido para solventar estos problemas y nos centramos en las técnicas basadas en funciones \textit{hash}, que permiten distribuir la carga de trabajo de forma más equilibrada y reducir la necesidad de sincronización global. En concreto, veremos una versión llamada HDA* (\textit{Hash Distributed A*}) que es uno de los resolvedores del estado del arte actual y que mejores resultados ha proporcionado en cuanto a su escalabilidad y eficiencia, además de ser una formulación paralela muy simple.

Como parte del trabajo, se han implementado tres versiones del algoritmo: una versión secuencial, una versión paralela centralizada (SPA*) y una versión paralela descentralizada basada en HDA*. Todas ellas se han desarrollado en el lenguaje C y utilizando OpenMP para gestionar el paralelismo en entornos de memoria compartida mediante hilos. Para la versión descentralizada se ha adaptado el diseño original de HDA* pensado para MPI, de forma que la comunicación entre procesos se simula mediante estructuras compartidas sincronizadas y el envío de nodos se hace de manera dinámica para evitar bloqueos innecesarios. Esta adaptación permite estudiar el comportamiento descentralizado sin necesidad de un entorno distribuido real.

El código se ha estructurado en módulos independientes para facilitar la comparación entre versiones, pudiéndose utilizar cualquiera de ellas desde un mismo ejecutable. La implementación es independiente de la estructura utilizada para representar el espacio de búsqueda. Se ofrece una interfaz para que el usuario defina las funciones de expansión y de cálculo de la heurística para adaptarlo a un dominio específico. Para una mayor eficiencia se ha utilizado una lista de nodos visitados en lugar de una lista de nodos cerrados, una técnica muy común en las implementaciones más modernas de A* y que nos permite prescindir de algunas operaciones sobre las estructuras de datos.

En cuanto a la experimentación, se han seleccionado una serie de escenarios de prueba con distintas características que permiten simular diferentes tipos de grafos: desde estructuras regulares con pocos obstáculos hasta mapas más irregulares con altos factores de ramificación. Estos mapas se han cargado a partir de ficheros .map y .scen ampliamente utilizados en benchmarks para este tipo de algoritmos y que están disponibles para su uso libre. Se han ejecutado las tres versiones del algoritmo variando el número de hilos utilizados en OpenMP en las versiones paralelas para observar su comportamiento en términos de tiempo de ejecución, escalabilidad y eficiencia.

Los resultados nos mostrarán que la versión secuencial funciona de forma robusta y es adecuada para problemas de tamaño moderado, pero su rendimiento cae rápidamente cuando el grafo crece en tamaño o complejidad. Observaremos rápidamente cómo la versión paralela centralizada obtiene en general peores resultados que la secuencial debido al cuello de botella introducido por la contención en los accesos a las estructuras de datos compartidas. Este resultado será coherente con lo descrito en la literatura, donde se observa que el \textit{speedup} obtenido con SPA* está acotado por el tiempo que los hilos pasan esperando para tener acceso a estructuras críticas. Sin embargo, comprobaremos que para problemas en los que la operación de expansión es costosa, sigue siendo una técnica muy viable a tener en cuenta, llegando a obtener un rendimiento incluso mejor que la versión descentralizada.

La versión descentralizada, por su parte, demostrará una mejor escalabilidad cuando se incrementen el número de hilos, aunque a costa de generar un mayor número de nodos y de introducir cierta sobrecarga en la búsqueda. Por otro lado, la necesidad de distribuir el trabajo entre los hilos también introduce una sobrecarga de comunicación. Veremos cómo estos problemas se hacen más notables en función del tipo de grafo y de sus características. Aun así, el diseño basado en funciones hash permite un reparto equilibrado del trabajo, y los resultados muestran mejoras significativas en tiempo cuando se utilizan escenarios grandes o espacios de búsqueda muy ramificados.

Las conclusiones del trabajo reflejarán cómo la estrategia descentralizada es la más adecuada en términos de escalabilidad y eficiencia, sobre todo para espacios de búsqueda grandes. Sin embargo, la versión centralizada conseguirá buenos resultados en entornos donde la operación de expansión es costosa. También analizaremos un parámetro referente a la versión descentralizada y observaremos cómo afecta el valor de éste al comportamiento del algoritmo, dejándonos ver que hay un valor óptimo. La posible automatización de la búsqueda de este valor óptimo se prpondrá como una posible línea de investigación en el futuro.

Finalmente, propondremos una serie de líneas futuras de investigación que pueden surgir a partir de este trabajo. Entre ellas se propondrá la posibilidad de realizar un análisis más a fondo de los algoritmos propuestos en términos del uso de memoria y generación de nodos. También se sugerirá implementar HDA* en un entorno de memoria distribuida según se expone en la publicación original y, por último, estudiar las implementaciones para arquitecturas de GPU, que aprovechan la alta paralelización de estas unidades de procesamiento para acelerar la expansión de nodos, pero que necesitan introducir técnicas efectivas de detección de duplicados en estos dispositivos.

\newpage

%%%%%%%%%%%%%%%%%%%%%%%%%%%%%%%%%%%%%%%%%%%%%%%%%%%%%%%%%%%%%%%%%%%%%%%%%%%%%%%
\section*{Extended Abstract}
%%%%%%%%%%%%%%%%%%%%%%%%%%%%%%%%%%%%%%%%%%%%%%%%%%%%%%%%%%%%%%%%%%%%%%%%%%%%%%%
\thispagestyle{empty}

In this work, we explore the different parallelization strategies for the A* algorithm in shared memory environments. The topic of this work was motivated by the completion of a Master’s degree in Quantum Computing at the \textit{Universidad Internacional de La Rioja} during the current academic year. In this Master’s program, a Master’s Thesis (TFM) line was proposed focusing on a quantum version of the A* algorithm. The main goal of this work has been to study the parallelization of the algorithm and its current state in classical computing, in order to establish a solid foundation for comparing it with a quantum proposal in the future.

The A* algorithm is a heuristic search algorithm used to solve the problem of finding the lowest-cost path between two nodes in a graph. This algorithm has proven to be a highly efficient technique, currently applied in a wide variety of fields such as video game development, intelligent transportation, robotic planning, and artificial intelligence. However, when attempting to formulate a parallel version, several complications arise—some inherent to parallel environments, and others due to the inherently sequential nature of the algorithm itself.

We begin this work with a review of how this problem has been addressed historically. We will examine when the first formulations and strategies to solve it emerged, and we will study it not only from the perspective of solving modern large-scale problems, but also as a problem for which nature has developed efficient adaptive mechanisms. We will present the mathematical foundations necessary to understand both the formulation of the shortest path problem and the operation of the A* algorithm itself. These include concepts from graph theory such as nodes, edges, directed and weighted graphs, as well as the notions of path, associated cost, and optimal path. We will formally state the formulation of the shortest path problem as the search for a sequence of nodes that minimizes the total cost from a start node to a goal node, under the condition that each consecutive node is connected via a valid edge.

Next, we briefly review the different search algorithm paradigms that exist today, classified into two main groups: informed and uninformed. Uninformed algorithms, such as Depth-First Search (DFS), Breadth-First Search (BFS), or Uniform Cost Search (UCS), lack information about the estimated distance to the goal, which leads them to expand a larger number of nodes and thus be less efficient. In contrast, informed or heuristic algorithms, such as Greedy Best-First Search (GBFS) or A*, use a heuristic function to guide the search toward the solution more efficiently.

The A* algorithm was proposed in 1968 by Hart, Nilsson, and Raphael in the context of the Shakey project, one of the first robots capable of reasoning about its actions. A*'s main contribution was the formal introduction of heuristic information into search problems. To do this, they introduced the use of this information to estimate the remaining cost to the goal and combined it with the accumulated cost from the start. The evaluation function they proposed for A*, which guides the search, is $f(n) = g(n) + h(n)$, where $g(n)$ represents the cost from the start to the current node and $h(n)$ is the estimate to the goal. This approach allows optimal paths to be found more efficiently than classical algorithms. We will see that, beyond being a single algorithm, A* actually encompasses a family of algorithms, with Uniform Cost Search and Greedy Best-First Search being particular cases. We present a detailed example of how it works to illustrate the process of node expansion and the use of the open and closed lists.

We will also discuss the numerous adaptations of the A* algorithm that have been developed to address domain-specific challenges or particular constraints, such as memory limitations and dynamic environments that may change during execution. Over the years, researchers have proposed various modifications to mitigate A’s most significant drawback—its substantial memory consumption, which grows exponentially with search depth in many cases. These adaptations include memory-efficient variants like IDA* (Iterative Deepening A*), which trades memory for computation by revisiting nodes, and SMA* (Memory-Bounded A*), which imposes strict memory constraints to prevent overflow. Other variants, such as D* (Dynamic A*) and its successors, focus on dynamic environments where the search space evolves in real time, allowing for efficient re-planning without restarting the search from scratch. Additionally, algorithms like Weighted A introduce heuristic weighting to balance optimality and speed, while approaches like Beam Search limit the number of nodes stored to control memory usage. These innovations highlight the flexibility of A*-based methods and their ability to be tailored to diverse applications, from robotics and pathfinding to real-time strategy games and logistics planning.

We will also analyze the fundamental properties of the algorithm, such as admissibility and optimality. We will see how these properties are related to the nature of the heuristic function used, revealing that the choice and design of the heuristic is one of the most important parts of the algorithm. The admissibility property guarantees that A* will find the optimal solution as long as the heuristic function does not overestimate the real cost (i.e., it is optimistic), while the optimality property refers to A* expanding the smallest number of nodes necessary to find the optimal solution. In particular, we show that if the heuristic is consistent (also called monotonic), the algorithm will not need to revisit the same node multiple times. For the study of optimality, we analyze the paper by Dechter and Pearl, where they propose different types of optimality for different kinds of algorithms depending on the characteristics of the problem being addressed. We conclude with a classification that specifies under which conditions and for what types of problems A* is an optimal algorithm.

Once we understand the sequential behavior of A*, we focus on the state of the art in its parallelization. Over the past decades, multiple strategies have been proposed to adapt this algorithm to parallel environments, especially driven by the need to reduce execution times for large-scale problems. These strategies can be broadly categorized into two approaches: centralized and decentralized parallelization.

In centralized strategies, all processes involved in the search share a global open list and a common closed list, which are the main data structures used by the algorithm. The most representative version of this approach is SPA* (Simple Parallel A), where processes compete to extract nodes from the open list and expand them, synchronizing through mutual exclusion mechanisms to avoid race conditions. This strategy is simple to implement in shared memory environments but presents significant scalability limitations. As the number of processes increases, contention for access to shared lists grows, and the wait time between processes can outweigh the gains from parallelization. Furthermore, the algorithm’s termination condition must be redefined to ensure that the best global solution has been found, as multiple processes may be expanding nodes in parallel without full knowledge of the search state. We will see how this type of strategy has been applied to domains where the expansion operation is costly and contention in shared data structures represents only a small part of the total algorithm execution time.

On the other hand, decentralized strategies distribute the work among processes such that each one maintains its own open and closed list. This strategy introduces new challenges, such as duplicate detection and the need to synchronize the search state among processes. However, it allows for greater scalability and improves performance in large-scale problems. We will examine different workload distribution strategies that have emerged to address these issues, focusing on techniques based on hash functions, which allow for more balanced workload distribution and reduced need for global synchronization. Specifically, we will look at a version called HDA* (Hash Distributed A), which is one of the current state-of-the-art solvers and has provided excellent results in terms of scalability and efficiency, in addition to being a very simple parallel formulation.

As part of this work, three versions of the algorithm have been implemented: a sequential version, a centralized parallel version (SPA*), and a decentralized parallel version based on HDA*. All of them have been developed in C using the OpenMP library to manage parallelism in shared memory environments via threads. For the decentralized version, the original HDA* design (intended for MPI) has been adapted so that communication between processes is simulated using synchronized shared structures, and node transmission is done dynamically to avoid unnecessary blocking. This adaptation allows the decentralized behavior to be studied without the need for a real distributed environment.

The code has been structured into independent modules to facilitate comparison between versions, allowing any of them to be used from a single executable. The implementation is independent of the structure used to represent the search space. An interface is provided so that the user can define the expansion and heuristic functions to adapt the algorithm to a specific domain. For greater efficiency, a visited-nodes list is used instead of closed nodes, a very common technique in more modern A* implementations that lets us avoid some operations on data structures.

For experimentation, a series of test scenarios with different characteristics were selected to simulate various types of graphs: from regular structures with few obstacles to more irregular maps with high branching factors. These maps were loaded from .map and .scen files, widely used in benchmarks for this type of algorithm and freely available for use. The three versions of the algorithm were executed, varying the number of threads used in OpenMP for the parallel versions, in order to observe their behavior in terms of runtime, scalability, and efficiency.

The results show that the sequential version runs robustly and is suitable for moderately sized problems, but its performance quickly deteriorates as the graph grows in size or complexity. We quickly observed that the centralized parallel version generally performs worse than the sequential version due to the bottleneck caused by contention in accessing shared data structures. This behavior is consistent with what is described in the literature, where it is noted that the speedup of SPA* is limited by the time threads spend waiting for access to critical structures. However, we found that for problems where the expansion operation is costly, it is a very viable technique to consider, even outperforming the decentralized version.

The decentralized version, in turn, demonstrates better scalability as the number of threads increases, although at the cost of generating more nodes and introducing some overhead in the search. Additionally, the need to distribute the work among threads also introduces synchronization overhead. We will see how these issues become more or less significant depending on the type of graph and its characteristics. Even so, the design based on hash functions enables balanced workload distribution, and the results show significant time improvements when large scenarios or highly branched search spaces are used.

The findings of this work demonstrate that the decentralized strategy outperforms centralized approaches in terms of scalability and efficiency, particularly when dealing with large search spaces and fast node generation rates. This advantage stems from its ability to distribute computational load effectively across multiple threads, reducing bottlenecks and improving overall performance.

However, the centralized version remains a viable option in scenarios where the node expansion operation is computationally expensive. In such cases, the centralized approach can achieve competitive results by leveraging sequential processing and minimizing overhead associated with task distribution.

Additionally, we examine a key parameter in our proposal of the HDA* algorithm related to the decentralized strategy, which governs the number of nodes a single thread can expand before redistributing them to their respective owner threads. Our analysis reveals that this parameter significantly influences the algorithm’s behavior, with performance peaking at an optimal value. This observation suggests that fine-tuning this parameter could further enhance efficiency, presenting a promising direction for future research in optimizing parallel heuristic search algorithms.

Finally, we propose several future research directions that may arise from this work. Among them is the possibility of conducting a more in-depth analysis of the proposed algorithms in terms of memory usage and node generation. Another proposal is to implement HDA* in a distributed memory environment as originally described in the literature, and finally, to explore implementations for GPU architectures, which leverage the high parallelization capabilities of these processing units to accelerate node expansion, but require effective duplicate detection techniques on these devices. After proposing these directions for future work, we will conclude with a brief overview of the contributions made in the field of quantum computing, an area that is currently emerging and being actively researched.

\newpage

\tableofcontents
\thispagestyle{empty}

\newpage

%%%%%%%%%%%%%%%%%%%%%%%%%%%%%%%%%%%%%%%%%%%%%%%%%%%%%%%%%%%%%%%%%%%%%%%%%%%%%%%
\section{Introducción} \label{sec:introduccion}
%%%%%%%%%%%%%%%%%%%%%%%%%%%%%%%%%%%%%%%%%%%%%%%%%%%%%%%%%%%%%%%%%%%%%%%%%%%%%%%

El Algoritmo A* es un algoritmo de búsqueda en grafos heurístico (o informado) que garantiza, bajo unas determinadas condiciones, encontrar el camino de coste mínimo entre un par de nodos. Este problema, conocido como el problema del camino más corto (\textit{Shortest Path Problem}), puede parecer trivial a simple vista, pero es fácil darse cuenta de que, a medida que se consideran un mayor número de nodos y aristas, la dificultad de resolverlo aumenta rápidamente. Podemos pensar que este tipo de problemas tienen una relevancia natural relacionada con las necesidades de los seres humanos, los animales o las plantas. Localizar la ruta más corta para, por ejemplo, conseguir comida o refugio, podría representar una ventaja en la adaptación de una especie a su entorno.

Lo cierto es que sí podemos encontrar en la naturaleza varios ejemplos de estrategias y acercamientos a este tipo de problemas. Uno de ellos es la forma en que las hormigas utilizan las feromonas para encontrar la ruta óptima entre la colonia y una fuente de alimento \cite{colorni1991distributed}. Tras explorar el espacio aleatoriamente, cuando una de las hormigas encuentra alimento, vuelve a la colonia de forma más o menos correcta, dejando un rastro de feromonas. Este rastro será seguido por otras hormigas, que lo reforzarán al volver si efectivamente encontraron alimento. Otro ejemplo son los túbulos que produce el moho \textit{Physarum polucephalum} para conectar distintas colonias con sus fuentes de alimentos. Con este último se realizó un experimento en el que, sobre una placa translúcida que replicaba los accidentes geográficos de Tokio, y colocando alimento en las ubicaciones de Tokio y sus ciudades periféricas, se observó que el moho construyó una red de túbulos muy cercana a la actual red ferroviaria de Tokio \cite{tero2010rules}. Este experimento les valió a sus investigadores el premio Ig Nobel de Planificación de Transporte en 2010.

Para los humanos no ha sido diferente. Con el auge de las civilizaciones, comienzan a surgir escenarios en los que la optimización de rutas es cada vez más relevante. En Mesopotamia y Egipto, la construcción de caminos y redes fluviales eficientes permitió el transporte de mercancías y personas. El Imperio Romano destacó por su vasta red de calzadas que conectaban todo su territorio, optimizando el comercio y la movilidad militar.

A pesar de todo esto, la investigación matemática formal en este problema comenzó relativamente tarde, sobre todo si la comparamos con la de otros problemas de optimización combinatoria. La primera referencia encontrada en la literatura sobre un problema similar la encontramos en \cite{wiener1873ueber}, donde se refiere en particular al problema de encontrar la salida en un laberinto. Se cree que esta aparición tardía puede deberse a que es un problema relativamente fácil de resolver, lo que además explica el hecho de que, cuando este problema comenzó a adquirir popularidad, varios investigadores independientes habían desarrollado métodos similares \cite{schrijver2012history}. Otra razón para explicar esta aparición tardía es el hecho de que hubo que esperar hasta el siglo XVIII para contar con un marco teórico matemático sobre el que trabajar este tipo de problemas. Se trata de la teoría de grafos, desarrollada por Leonhard Euler, y que utilizó para resolver el famoso problema de los puentes de Königsberg \cite{euler1741solutio}.

Desde entonces, la teoría de grafos ha evolucionado significativamente y, junto con los avances en ciencias de la computación, ha motivado la investigación en el área de los algoritmos de búsqueda. Además, con ayuda de conceptos como el de la representación atómica para el modelado de problemas \cite[p. 57]{russell2016artificial}, el ámbito de aplicación se ha extendido mucho más allá de la búsqueda de un camino más corto, siendo de especial interés en áreas como la inteligencia artificial, la bioinformática, la toma de decisiones o la optimización de redes de transporte.

\subsection{Conceptos matemáticos} \label{sec:conceptos}

En esta sección vamos a introducir algunas definiciones matemáticas que nos ayudarán a entender el formalismo sobre el que se sustentan los métodos de búsqueda en grafos, así como una descripción del problema del camino más corto y los conceptos esenciales para comprender el funcionamiento del algoritmo.

\begin{defn}\label{def1}
Un \textbf{grafo} $G$ es un par ordenado $G=(V,E)$, donde $V$ es un conjunto de elementos $\{n_i\}$ llamados vértices (o nodos), y $E$ es un conjunto de elementos $\{e_{ij}\}$ llamados aristas (o arcos). Si el elemento $e_{pq}$ pertenece al conjunto $\{e_{ij}\}$, entonces decimos que existe un arco del nodo $n_p$ al nodo $n_q$.
\end{defn}

\begin{figure}[H]
    \centering
    \begin{minipage}{0.47\textwidth}
        \centering
        \includesvg[inkscapelatex=false, width=0.5\linewidth]{images/undirected.svg}
        \caption{Grafo no dirigido}
        \label{fig:undirected}
    \end{minipage}
    \begin{minipage}{0.47\textwidth}
        \centering
        \includesvg[inkscapelatex=false, width=0.5\linewidth]{images/directed.svg}
        \caption{Grafo dirigido}
        \label{fig:directed}
    \end{minipage}
\end{figure}

\noindent Podemos hablar de \textbf{grafo no dirigido} (Figura \ref{fig:undirected}) cuando los elementos de $E$ son pares no ordenados, de manera que $e_{pq}=e_{qp}$, y de \textbf{grafo dirigido} (Figura \ref{fig:directed}) cuando se trata de pares ordenados (en general, $e_{pq} \neq e_{qp}$). En estos últimos, la existencia del elemento $e_{pq}$ no implica la existencia del elemento $e_{qp}$. Para poder realizar una buena representación del problema del camino más corto, trabajaremos con \textbf{grafos dirigidos ponderados} (Figura \ref{fig:weighted}).

\begin{defn}
Un \textbf{grafo dirigido ponderado} es un grafo dirigido $G=(V,E)$, en el que cada elemento $e_{ij}\in E$ tiene asociado un coste $c_{ij} \in \mathbb{R}$.
\end{defn}

\begin{figure}
    \centering
    \includesvg[inkscapelatex=false, width=0.26\linewidth]{images/weighted.svg}
    \caption{Grafo dirigido ponderado.}
    \label{fig:weighted}
\end{figure}

\noindent Consideraremos solo aquellos grafos dirigidos ponderados para los que el coste de cualquier arista es positivo y no nulo.

Esta forma explícita de definir un grafo a partir de sus conjuntos de nodos y aristas atiende a la manera convencional en teoría de grafos, pero cuando se trata con estos problemas desde el punto de vista de la computación, es usual utilizar otra forma equivalente. Se trata de una definición implícita a partir de un conjunto de nodos fuente $S\subset V$ y un operador de vecindad $\Gamma$ definido sobre el conjunto de elementos $\{n_i\}$ cuyo valor para cada $n_i$ es un conjunto de pares $\{(n_j, c_{ij})\}$. De esta forma, a partir de un nodo inicial y un operador de vecindad podemos obtener el grafo completo.

\begin{defn}
Un \textbf{camino} de $n_1$ a $n_k$ es un conjunto ordenado de nodos $P = (n_1,n_2, \allowbreak \dots ,n_k)$ donde cada nodo $n_{i+1}$ es sucesor del nodo $n_i$. En un grafo ponderado, el camino entre tendrá un coste asociado $C(P) = \sum_{i=0}^{k-1}c_{i,i+1}$.
\end{defn}

\noindent 

\begin{defn}
El \textbf{camino óptimo} del nodo $n_i$ al nodo $n_k$ es el camino con el menor coste de entre todos los posibles caminos de $n_i$ a $n_k$. A este coste lo representamos como $h(n_i, n_k)$.
\end{defn}

\noindent Por simplicidad, y cuando el nodo objetivo $n_k$ es único, podemos representar el coste del camino óptimo desde un nodo $n_i$ al nodo objetivo como $h(n_i)$. Los algoritmos heurísticos hacen uso de esta función, denominada \textbf{heurística} para discriminar cuál de las siguientes opciones es más prometedora (eligiendo aquella que minimiza o maximiza su valor). Diremos que un algoritmo es \textbf{admisible} si está garantizado que encontrará la solución óptima.

Ahora, ya podemos formular el problema a resolver y explicar el funcionamiento del algoritmo.

\begin{defn}\label{def:shortest_path_problem}
Dado un grafo dirigido con pesos positivos y dos nodos $s$ y $t$, el \textbf{problema del camino más corto} consiste en encontrar una secuencia de nodos $P = (n_1, n_2 \dots, n_k)$ tal que $n_0 = s$ y $n_k = t$, cada par consecutivo $n_i, n_{i+1}$ está conectado por una arista del grafo, y el coste total del camino es el mínimo posible de entre todos los caminos que conectan $s$ con $t$.
\end{defn}

%{\color{blue} lo anterior lo pondría o bien indexado o bien redactado, pero no mezclado como está ahora}


\subsection{Algoritmos de búsqueda}\label{sec:alg-busqueda}

Un algoritmo de búsqueda es un procedimiento sistemático que explora un espacio de posibles soluciones con el objetivo de encontrar una o varias que cumplan ciertas condiciones, típicamente una solución óptima o satisfactoria, a un problema. Estos algoritmos están estrechamente relacionados con el problema del camino más corto, pues son la base para las estrategias que se llevan a cabo en este problema. 

A día de hoy podemos clasificar los algoritmos de búsqueda en dos categorías, los no informados y los informados o heurísticos. La diferencia principal está en que los algoritmos no informados no poseen ningún tipo de conocimiento sobre el dominio del problema, mientras que los informados disponen de alguna forma de evaluar si una opción es más prometedora que otra. Esta limitación de los algoritmos no informados con respecto a los informados hace que, en general, su eficiencia sea mucho menor.

Las distintas estrategias para realizar búsquedas no informadas se basan en explorar el espacio de búsqueda de alguna manera más o menos ordenada, hasta dar con el objetivo, sin tener más información que el poder saber si el estado actual es el objetivo o no. Las dos principales estrategias son la búsqueda en profundidad (DFS) y la búsqueda en anchura (BFS), que se diferencian en el orden en que se exploran los estados vecinos. El objetivo de estos algoritmos es el de encontrar si existe dicho estado objetivo o no, pero no resuelven el problema del camino más corto. Para ello, debemos tener en cuenta el coste acumulado desde el inicio hasta el nodo objetivo. Ejemplos de estos algoritmos son la búsqueda de coste uniforme (UCS) o el algoritmo de Dijkstra \cite{dijkstra_note_1959}.

Los algoritmos de búsqueda informados se basan en el diseño de una función heurística. Esta función permite evaluar los distintos estados del espacio de búsqueda para obtener alguna medida sobre lo prometedores que son, en cuanto a encontrar el objetivo se refiere. Uno de los principales focos de investigación en este área es precisamente el diseño de estas funciones heurísticas, y puede demostrarse para ciertos problemas que, si esta función cumple unas determinadas propiedades, el algoritmo encontrará la solución óptima. Un ejemplo de algoritmo de búsqueda heurístico es el algoritmo primero el mejor (GBFS), que selecciona en cada paso el nodo con el valor heurístico más bajo, es decir, el que parece estar más cerca del objetivo según la estimación. Sin embargo, este algoritmo no garantiza optimalidad, ya que puede quedar atrapado en soluciones subóptimas debido a que ignora el coste acumulado desde el nodo inicial.

El algoritmo A* fue propuesto en 1968 por Peter Hart, Nils Nilsson y Bertram Raphael \cite{hart_formal_1968} en el contexto del proyecto Shakey, un robot pionero en el campo de la inteligencia artificial con capacidad para razonar sobre sus propias acciones. En esta publicación introducen formalmente el uso de la información heurística sobre el dominio de un problema en las estrategias de búsqueda, y proporcionan la definición de los conceptos de admisibilidad y optimalidad. Se dice que un algoritmo de búsqueda es admisible cuando garantiza encontrar la solución óptima, y es óptimo cuando la encuentra explorando el mínimo número de estados necesarios para llegar a dicha solución. Bajo este marco, demuestran la admisibilidad y optimalidad del algoritmo A* en términos de dos propiedades que debe cumplir la función heurística utilizada: la admisibilidad y la consistencia. Decimos que una heurística es admisible si nunca sobrestima el coste real, y que es consistente si el valor heurístico de cualquier estado es menor que el coste de moverse a un estado vecino más el valor heurístico del vecino si el estado solución es más cercano al vecino que al estado actual. En las Secciones \ref{sec:admisibilidad} y \ref{sec:optimalidad} demostramos la admisibilidad y la optimalidad del algoritmo, y veremos cómo surgen naturalmente estas propiedades de la heurística.

\subsection{El algoritmo A*} \label{sec:algoritmo}

El algoritmo A* (A estrella) es un algoritmo de búsqueda informada que combina las ventajas de la búsqueda de coste uniforme y la búsqueda primero el mejor (velocidad mediante heurísticas). Su objetivo es encontrar el camino de menor coste desde un nodo inicial, $s$, hasta un nodo objetivo, $t$, en un grafo dirigido ponderado. El elemento diferenciador de este algoritmo es su función de evaluación, que se define como:

\begin{equation}
f(n) = g(n) + h(n)
\end{equation}

\noindent donde $g(n)$ representa el coste real acumulado desde un nodo inicial hasta el nodo $n$ por un camino óptimo, y $h(n)$ el coste real para llegar al nodo objetivo desde el nodo $n$ por un camino óptimo. Si solo consideramos la función $g(n)$, el algoritmo A* sería equivalente al algoritmo de búsqueda de coste uniforme y, si solo consideramos $h(n)$, entonces es equivalente al algoritmo primero el mejor. El hecho de tener en cuenta ambas funciones hace que el algoritmo nos permita no solo encontrar la solución óptima, sino también de manera más eficiente.

El objetivo del algoritmo es encontrar el camino que minimiza el valor de $f$, pero no siempre podemos conocer con exactitud los valores de $h$ y $g$, por lo que tendremos que utilizar una estimación de estos, es decir:

\begin{equation}
\hat{f}(n) = \hat{g}(n) + \hat{h}(n)
\end{equation}

La propuesta de Hart et al. \cite{hart_formal_1968} para la elección de $\hat{g}(n)$ es el menor coste encontrado hasta el momento actual entre el nodo inicial y $n$. Esta elección implica que $\hat{g}(n) \geq g(n)$. La elección de $\hat{h}(n)$ es más complicada y requiere de un conocimiento explícito del dominio del problema. En algunos casos, como por ejemplo, encontrar el camino entre dos puntos de un mapa, la distancia en línea recta sirve como estimación, pero en otros dominios puede no ser una elección trivial. De hecho, uno de los principales problemas que se presentan al querer aplicar este algoritmo en un problema real es determinar esta estimación de $h$.

La manera de proceder del algoritmo consiste en aplicar el operador de vecindad al estado actual para obtener sus estados vecinos y explorar aquel con un valor de $\hat{f}$ más pequeño. Para tener en cuenta aquellos estados que ya han sido visitados, y los que quedan por visitar, se utilizan dos estructuras de datos llamadas lista abierta, $Q$, y lista cerrada $H$ \footnote{Las listas abierta y cerrada se llaman así por razones históricas. De hecho, no suelen implementarse como listas en la práctica.}. La lista abierta almacena aquellos estados a los que podemos llegar pero aún no hemos explorado, y la lista cerrada aquellos que ya han sido explorados y, además, sus vecinos se han introducido en la lista abierta. Decimos que un estado está \textit{abierto} o \textit{cerrado} si está en la lista abierta o cerrada, respectivamente.

El algoritmo, tal como se define en \cite{hart_formal_1968}, opera de la siguiente manera:

\begin{enumerate}
\item Marcar $s$ como abierto y calcular $\hat{f}(s)$. \label{as:p1}
\item Seleccionar el nodo abierto $n$ con el mínimo valor de $\hat{f}(n)$. \label{as:p2}
\item Si $n = t$, marcar $n$ como cerrado y terminar el algoritmo. \label{as:p3}
\item Si no, marcar $n$ como cerrado y aplicar el operador $\Gamma$ a $n$. Calcular $\hat{f}(n)$ para cada vecino de $n$ y marcar como abiertos aquellos que no están marcados como cerrados. Remarcar como abiertos los sucesores ya marcados cerrados $n_i$ para los que el valor de $\hat{f}(n_i)$ es menor ahora que cuando fueron marcados cerrados. \label{as:p4}
\end{enumerate}

\noindent A la hora de implementar el algoritmo, queda claro que son necesarias dos estructuras de datos principales, la lista abierta y cerrada. Para determinar la mejor implementación se deben revisar las principales operaciones que se realizan sobre ellas. Comenzando por la lista cerrada, al ser la encargada de almacenar los nodos expandidos, realizaremos la operación de inserción y, al expandir un nodo, debemos comprobar si cualquier nodo adyacente está en la lista, por lo que haremos una consulta (Paso \ref{as:p4}). La estructura de datos más común es la tabla hash, que soporta ambas operaciones en tiempo constante en el mejor de los casos, esto es, cuando la función hash utilizada no provoca colisiones. El caso más sencillo y que podemos encontrar en algunas implementaciones es utilizar un array e identificar a cada nodo de forma unívoca por un número entero, que será su posición en el array. En cuanto a la lista abierta, es necesario extraer de ella el nodo con el menor valor $\hat{f}$, y también habrá que insertar los nuevos nodos que aún no han sido expandidos (Paso \ref{as:p4}). El hecho de que se extraigan los nodos de la lista abierta por orden de prioridad de $\hat{f}$ hace que la estructura de datos ideal sea una cola de prioridad. Según cómo se implemente la cola de prioridad, se obtienen órdenes de complejidad diferentes para las operaciones de inserción y extracción. Lo más habitual es implementar esta lista como un montículo binario, que permite realizar ambas operaciones en tiempo $O(\log n)$. Un montículo binario es un array en el que la información contenida mantiene una estructura de tipo árbol binario, de manera que cada nodo tiene un valor de prioridad menor (o mayor) que sus dos hijos. La forma de estructurar el array es la siguiente: cada nodo en la posición $n$ tendrá a sus hijos en las posiciones $2n$ y $2n+1$ si empezamos a contar las posiciones desde el uno, y en las posiciones $2n + 1$ y $2n + 2$ si empezamos a contar desde el cero (Figura \ref{fig:binheap}).

\begin{figure}[H]
    \centering
    \includesvg[scale=0.8]{images/heap-tree.svg}
    \caption{Representación en forma de árbol y array de un montículo binario. El valor representado en cada nodo indica su prioridad (coste $\hat{f}$) y las flechas apuntan a los hijos de cada nodo.}
    \label{fig:binheap}
\end{figure}

\noindent Aunque lo más habitual es implementar la cola de prioridad como un montículo binario, también se pueden encontrar otras estructuras de datos como montículos de Fibonacci \cite{fredman1987fibonacci}. Éstos, en particular, son útiles porque permite reubicar un nodo en el array en tiempo constante cuando se modifica su prioridad. A partir de ahora nos referiremos a la lista abierta también como la cola de prioridad de manera indistinta.

A continuación podemos ver el pseudocódigo del algoritmo. Los valores de $\hat{g}(n)$, $\hat{f}(n)$ y $parent(n)$ son los debemos almacenar para cada nodo. Las implementaciones más sencillas utilizan tres arrays distintos y prescinden de la lista cerrada como estructura de datos en sí misma, pero lo más habitual es crear alguna estructura de datos que represente a los nodos, y almacenar en la lista cerrada los nodos con esos valores.

\begin{algorithm}[h]
\caption{A*}\label{astar}
\begin{algorithmic}[1]
\State $Q \gets \{s\}$;
\State $H \gets \emptyset$;
\While{$Q \neq \emptyset$}
    \State $n \gets \textsc{ExtractMin}(Q)$;
    \State $\textsc{Insert}(H, n)$;
    \If{$n = t$}
        \State \textbf{return} \textsc{RetracePath}($n$);
    \EndIf
    \ForAll{$n' \in \textsc{Neighbors}(n)$}
        \State $g' \gets \hat{g}(n)+c(n, n')$;
        \If{$n' \in H$}
            \If{$g' < \hat{g}(n')$}
                \State $\textsc{Remove}(H, n')$;
                \State $\textsc{Insert}(Q, n')$;
            \Else
                \State \textbf{continue};
            \EndIf
        \Else
            \If{$n' \notin Q$}
                \State $\textsc{Insert}(Q, n')$;
            \ElsIf{$g' \geq \hat{g}(n')$}
                \State \textbf{continue};
            \EndIf
        \EndIf
        \State $\hat{g}(n') \gets g'$;
        \State $\hat{f}(n') \gets g(n') + h(n')$;
        \State $parent(n') \gets n$;
    \EndFor
\EndWhile
\State \textbf{return} error $Q$ is empty;
\end{algorithmic}
\end{algorithm}

El algoritmo comienza creando la lista abierta, que incluye al nodo de inicio, y la lista cerrada, inicialmente vacía (líneas 1-2). Acto seguido entra en el bucle principal (líneas 3-29), que ejecutará mientras la lista abierta contenga elementos. Se extrae el nodo $n \in Q$ con menor coste $\hat{f}$ y se inserta en la lista cerrada (líneas 4-5). Si el nodo extraído es el objetivo, el algoritmo reconstruye el camino desde $n$ y termina su ejecución (líneas 6-8). En caso de no ser el objetivo, se expande el nodo y se consideran sus vecinos (líneas 9-28). Se calcula el nuevo coste para alcanzar al vecino a partir del nodo actual $g'$ (línea 10) y se comprueba si ya ha sido cerrado. En ese caso, si el nuevo coste $g'$ mejora al coste anterior $\hat{g}(n')$, el nodo debe ser reabierto, esto es: se elimina de la lista cerrada, se inserta en la lista abierta (líneas 13-14) y se actualizan sus valores (líneas 15-17). En caso de no mejorar, podemos descartar el nodo y seguir (línea 16). Si el nodo vecino no ha sido cerrado, debemos comprobar si está o no en la lista abierta pendiente de expandir (línea 19). Si no está, se trata de un nuevo nodo, por lo que lo insertamos (línea 20) y actualizamos sus valores (líneas 25-27). Si está en lista abierta y el nuevo coste no mejora al anterior, podemos descartarlo (línea 22).


Cabe destacar que cuando la heurística es consistente, ningún nodo será reabierto, pues esta propiedad nos garantiza que cuando un nodo es cerrado, el coste óptimo para llegar a él ha sido encontrado, $\hat{g}(n) = g(n)$. Podemos encontrar otras versiones del pseudocódigo en las que se sustituye esta sección por un \textit{continue} (líneas 12-17) cuando esta condición de consistencia está garantizada.

Al ser un algoritmo heurístico, la eficiencia de éste dependerá mucho de las característica de la función elegida como heurística. Cuanto más parecida sea la estimación al valor real de $h$, menos nodos expandirá el algoritmo. Sabemos que, como mucho, A* expandirá todos aquellos nodos con un valor $f(n) \leq C^*$, siendo $C^*$ el coste del camino óptimo. Dorian y Michie \cite{doran_experiments_1997} proponen una medida llamada \textit{penetrancia}, $P$, que se define como la longitud del camino óptimo, $L$, entre el número total de nodos generados durante la búsqueda, $T$, incluyendo al nodo objetivo pero no al inicial.

\begin{equation}
P = \frac{L}{T}
\end{equation}

\noindent Cuando este valor es cercano a uno, quiere decir que el algoritmo ha expandido pocos nodos que no están en el camino óptimo y, cuando es cercana a cero, nos encontraremos que el algoritmo se comporta casi como una búsqueda no informada. Este valor depende tanto de la dificultad del problema como de la eficiencia del método de búsqueda. En la literatura más moderna en relación con la inteligencia artificial, se puede ver la eficiencia del algoritmo expresada en el peor caso como $O(b^d)$, donde $b$ es el factor de ramificación y $d$ la profundidad del nodo objetivo. En cuanto a la complejidad en espacio, A* guarda todos los nodos generados en la lista cerrada y los nodos por explorar en la lista abierta, y en cada iteración visita un nodo, por lo que la complejidad en espacio es también exponencial con respecto a la longitud del camino $d$.

En la práctica, éste es posiblemente el mayor problema que tiene el algoritmo, y ha llevado a desarrollar versiones limitadas en memoria para dominios de problema en los que el espacio de búsqueda crece de forma exponencial. Algunos de estos algoritmos son: IDA* \cite{korf1985depth}, que realiza una búsqueda hasta una profundidad limitada y que incrementa iterativamente si no encuentra la solución; SMA* \cite{russell1992cient}, que elimina los nodos menos prometedres cuando la memoria está cerca de llenarse; RA* \cite{evett1995massively}, que utiliza una operación llamada \textit{retracción} para almacenar la información de los nodos de la frontera en sus padres, ahorrando memoria. Además de estas optimizaciones, como la eficiencia del algoritmo está íntimamente relacionada con la naturaleza del problema, se han creado a lo largo de los años otras versiones enfocadas a dominios de problema específicos: Anytime A* \cite{hansen2007anytime}, Any-angle A* \cite{daniel2010theta}, D* \cite{stentz1994optimal} o LPA* \cite{koenig2001incremental}.

Esta gran variedad de \textit{sabores} diferentes del algoritmo a dado lugar a la aplicación y perfección de éste en dominios específicos. En el ámbito del desarrollo de videojuegos, es habitual representar el espacio de búsqueda como una rejilla \cite{goos_grid-based_2002, harabor_online_2011, standley_complete_2011, barnouti_pathfinding_2016} y, además, a veces el camino debe buscarse para varios personajes \cite{felner_adding_2018}. Una de las variedades de A* que más se aplica en los videojuegos a día de hoy es LRTA* \cite{bulitko2006learning} ya que se trata de un algoritmo en tiempo real y que utiliza espacios de búsqueda locales en lugar del grafo completo para ahorrar memoria. En el artículo de \cite{rafiq_pathfinding_2020} podemos ver como las estrategias basadas en heurísticas han ido sobreponiéndose a las basadas en algoritmos genéticos en los videojuegos. También en el campo de la robótica son útiles estas técnicas para resolver problemas como la planificación de rutas \cite{martins_improved_2022}.

\subsubsection{Ejemplo de uso}

Con el fin de comprender en detalle el funcionamiento del algoritmo, más allá e una descripción de éste, vamos a ver un ejemplo de uso. Consideremos el grafo de la Figura \ref{fig:example}, donde el nodo $A$ es el nodo inicial y el nodo $F$ es el objetivo:

\begin{figure}[H]
\centering
\includesvg[width=0.6\textwidth]{images/example.svg}
\caption{Grafo de ejemplo. Los valores en negro respresentan los costes de recorres cada arista y los valores en gris representan la estimación de la heurística para cada nodo.}
\label{fig:example}
\end{figure}

Una vez inicializadas la lista abierta y la lista cerrada, el algoritmo comienza evaluando el nodo A e introduciéndolo en la lista abierta, $S = \{A\}$, con un valor $\hat{f}(A) = \hat{g}(A) + \hat{h}(A) = 0 + 11 = 11$. Acto seguido, entra en el bucle principal, en el que permanecerá mientras no se encuentre una solución y $S \neq \emptyset$. Se extrae $A$ de la lista abieta y se introduce en la lista cerrada. A continuación se expande el nodo $A$ y se obtiene como vecinos los nodos $B$ y $C$ con un coste para alcanzarlos de 5 y 4, respectivamente. Como son nodos nuevos, se evalúa su coste $\hat{f}$ y insertan en la lista abierta. En este caso: $\hat{f}(B) = 5 + 8 = 13$ y $\hat{f}(C) = 4 + 8 = 12$, y nos quedan $Q = \{B, C\}$ y $H=\{A\}$. En la siguiente iteración se extraerá de la lista abierta el nodo $C$, por tener el menor valor de $\hat{f}$. Se introduce $C$ en la lista cerrada y se evalúan sus vecinos, que en este caso son los nodos $B, D$ y $E$. Para el nodo $B$ obtenemos un nuevo coste $\hat{f}'(B) = 5 + 8 = 13$ ($4 + 1 = 5$ es el coste de alcanzar $B$ pasando por $C$), como $B \in Q$ y este nuevo coste no mejora al anterior, podemos descartar este camino. Los nodos $D$ y $E$ son nuevos, por lo que se evalúan: $\hat{f}(D) = 12 + 4 = 16$ y $\hat{f}(E) = 14 + 2 = 16$ y se insertan en la lista abierta, quedando $Q = \{B, D, E\}$ y $H = \{A, C\}$. El próximo nodo que se extraerá de la lista abierta será $B$, con $\hat{f}(B)= 13$, y se inserta en la lista cerrada. Desde $B$ podemos llegar a $C$ con un coste $\hat{f}'(C) = 6 + 8 = 14$, que se descarta por no mejorar el valor acutal; y a $D$, con coste $\hat{f}'(D) = 10 + 4 = 14$. En el caso de $D$ hemos mejorado el coste anterior, por lo que actalizamos a este nuevo valor y guardamos el camino para llegar a $D$ desde $B$. Nos queda $Q = \{D, E\}$ y $H = \{A, B, C\}$. A continuación se extrae $D$ de la lista abierta con coste $\hat{f}(D) = 14$ y se introduce en la lista cerrada. Desde $D$ se puede alcanzar $E$ con coste $\hat{f}'(E) = 12 + 2 = 14$, que mejora al coste anterior, por lo que actualizamos el camino hacia $E$; y también podemos alcanzar $F$ con coste $\hat{f}(F) = 16$, que introducimos en la lista abierta al ser un nuevo nodo. Nos queda $Q = \{E, F\}$ y $H = \{A, B, C, D\}$. En la siguiente iteración se extraería $E$ de la lista abierta, y al evaluar sus vecinos no se actualizaría ninguno por no mejorar los caminos ya encontrados. Finalmente, se extraerá $F$, y al tratarse de un nodo objetivo, podemos terminar la búsqueda y devolver el camino óptimo, en este caso: $A \rightarrow B \rightarrow D \rightarrow E \rightarrow F$.

\subsubsection{Admisibilidad de A*}\label{sec:admisibilidad}

En esta Sección vamos a discutir la admisibilidad del algoritmo, es decir, probaremos que dado un problema, encuentra la solución óptima. Nos basaremos en la prueba dada por uno de los autores del algoritmo en un libro publicado posteriormente \cite[p. 59]{nillsson1971problem}. En la Sección \ref{sec:alg-busqueda} introducíamos el concepto de algoritmo admisible como aquel que garantiza encontrar la solución óptima. Esta propiedad está directamente ligada a las características de la función heurística y a la elección de la estimación $\hat{h}$. Más concretamente, probaremos que el algoritmo A* es admisible si $\hat{h}(n) \leq h(n)$ para todo $n$.

\begin{lem}\label{lem}
Si $\hat{h}(n) \leq h(n)$ para todo $n$, entonces en cualquier momento antes de que A* finalice y para cualquier camino $P$ desde el nodo $s$ hasta el nodo $t$, existe un nodo abierto $n'$ en $P$ con $\hat{f}(n') \leq f(s)$.

\end{lem}

\begin{proof}
Sea un camino óptimo $P = (s=n_0, n_1, n_2, \dots, n_k = t)$. Debe haber al menos un nodo $n' \in P$ abierto, pues si $n_k$ fuera cerrado, el algoritmo habría terminado. Por definición de $\hat{f}$ tenemos:

$$\hat{f}(n) = \hat{g}(n) + \hat{h}(n)$$

\noindent Sabemos que A* ha encontrado el camino óptimo a $n'$, pues $n'\in P$ y todos sus antecesores en $P$ están cerrados. Por tanto $\hat{g}(n') = g(n')$ y:

$$\hat{f}(n') = g(n') + \hat{h}(n')$$

\noindent Como hemos asumido $\hat{h}(n) \leq h(n)$ podemos escribir:

$$\hat{f}(n') \leq g(n') + h(n') = f(n')$$

\noindent Pero el valor $f$ de cada nodo en un camino óptimo es igual a $f(s)$, el coste mínimo, y por tanto $\hat{f}(n') \leq f(s)$.
\end{proof}

Podemos probar ahora que el algoritmo es admisible si se cumple que $\hat{h}$ es una estimación optimista de $h$.

\begin{thm}\label{thm1}
Si $\hat{h}(n) \leq h(n)$ para todos los nodos $n$, y si todos los costes de los arcos son mayores o iguales que un número positivo $\delta$, entonces el algoritmo $A*$ es admisible.
\end{thm}

\begin{proof}
Probaremos este teorema por reducción al absurdo. Asumiremos que el algoritmo termina sin encontrar el camino óptimo. Entonces hay tres casos a considerar: termina sin encontrar el nodo objetivo, no termina nunca o termina pero no encuentra el camino de coste mínimo.

\begin{case1}
    Termina sin encontrar el nodo objetivo. Este caso contradice la condición de terminación del algoritmo (Paso \ref{as:p3}), por lo que no se considera. El único otro caso en el que puede terminar el algoritmo se da cuando la lista abierta está vacía, pero por el Lemma \ref{lem} sabemos que esta situación no puede darse si existe un camino entre el nodo inicial y el final.
\end{case1}

\begin{case2}
    No termina nunca. Supongamos un nodo objetivo $t$ accesible desde un nodo inicial $s$ en un número finito de pasos y con coste asociado $f(s)$. Dado que el valor de cada arco es positivo y mayor que cero, entonces para cada nodo $n$ a $M=f(s)/\delta$ pasos de $s$ se tiene $\hat{f}(n) \geq \hat{g}(n) \geq g(n) > M\delta = f(s)$. De esta expresión podemos ver que ningún nodo $n$ a más de $M$ pasos del nodo inicial $s$ será expandido nunca, dado que por el Lemma \ref{lem}, habrá un nodo $n'$ en la lista abierta con $\hat{f}(n') \leq f(s) < \hat{f}(n)$, y A* elegirá a $n'$ antes que a $n$. Aun así, el A* podría no terminar debido a reabrir continuamente nodos a menos de $M$ pasos de $s$. Sea $\chi(M)$ el conjunto de nodos accesibles desde $s$ en $M$ pasos, y sea $\nu(M)$ el número de nodos en $\chi(M)$. Cada nodo $n\in\chi(M)$ podrá ser reabierto a lo sumo un número finito de veces $\tilde{\rho}(n, M)$, dado que hay un número finito de caminos de $s$ a $n$ pasando solamente por nodos a $M$ pasos de $S$. Sea
    $$\rho(M) = \max_{n\in\chi(M)}\tilde{\rho}(n, M)$$
    El máximo número de veces que un nodo puede reabrirse. Por tanto, tras como mucho $\nu(M)\rho(M)$ expansiones, todos los nodos en $\chi(M)$ habrán sido cerrados. Dado que ningún nodo fuera de $\chi(M)$ será nunca expandido, A* termina.
\end{case2}

\begin{case3}
    Termina pero no encuentra el camino de coste mínimo. Supongamos que A* termina en algún nodo objetivo $t$ con $\hat{f}(t)=\hat{g}(t)> f(s)$. Por el Lemma \ref{lem}, habrá un nodo en la lista abierta $n'$ con $\hat{f}(n) \leq f(s) < \hat{f}(t)$, por lo que A* elegirá a $n'$ antes que a $t$, contradiciendo que el algoritmo terminaría en un camino no óptimo.
\end{case3}

\end{proof}

La prueba del Teorema \ref{thm1} queda terminada. Hemos completado la prueba del Teorema \ref{thm1} y podido comprobar, como anunciábamos al comienzo de la sección, que cuando la estimación de la heurística es un límite inferior de ésta, $\hat{h}(n) \leq h(n)$, para todo $n$, entonces A* es admisible, y encuentra siempre la solución óptima.

\subsubsection{Optimalidad de A*}\label{sec:optimalidad}

La optimalidad de A* fue una cuestión que condujo a confusión, y hubo que esperar al desarrollo de un marco de trabajo más completo para algoritmos de búsqueda informada. Como se ha mencionado anteriormente, decimos que un algoritmo de búsqueda es óptimo cuando expande el mínimo número de nodos necesarios para encontrar la solución óptima. La primera prueba de la optimalidad de A* la ofrecieron Hart et al. \cite{hart_formal_1968}, que posteriormente sintetizaron en el libro de Nilsson \cite{nillsson1971problem}. En su prueba, imponían la condición de consistencia a la función heurística. Esto es, dado un nodo $n$ y un nodo $n'$ tal que $h(n) < h(n')$, entonces debe cumplirse:

\begin{equation}
\hat{h}(n) < c(n,n') + \hat{h}(n')
\end{equation}

\noindent Cabe destacar que la propiedad de consistencia implica la propiedad de admisibilidad, pero en la dirección contraria no es cierto. Años más tarde, publicaron una corrección \cite{hart_correction_1972} en la que sostenían que la condición de consistencia no era necesaria, pues no hacían uso de ella en su demostración, pero fueron refutados por Dechter y Pearl \cite{dechter1985generalized, dechter1983optimality}. Se basaron en la observación de Gelperin \cite{gelperin1977optimality}, quien había notado que para discutir sobre la optimalidad del algoritmo había que compararlo con una clase más amplia de algoritmos igualmente informados, no solo aquellos que hacían uso de la heurística en la forma $f = g + h$. Dechter y Pearl se dieron cuenta de que la definición de Gelperin de \textit{igualmente informado} no incluía a aquellos algoritmos que, teniendo la misma información disponible que A*, pudieran hacer un mejor uso de ésta. El resultado de su trabajo es una jerarquía de cuatro tipos de optimalidad para tres tipos de algoritmos y cuatro dominios de problemas. No entraremos en los detalles de las demostraciones por la complejidad que presentan en comparación con la prueba de admisibilidad de algoritmo, pero vamos a repasar los resultados Dechter y Pearl, así como las conclusiones más notables.

Una instancia de un problema se define como una cuádrupla $I=(G, s, \Gamma, \hat{h})$, donde $G$ es el grafo, $s$ el nodo inicial, $\Gamma$ el operador de vecindad y $\hat{h}$ las estimaciones de la heurística. Dechter y Pearl proponen los siguientes conjuntos de instancias relacionados con la admisibilidad y la consistencia de la heurística:

\begin{equation}
\begin{split}
& \textbf{I}_{AD} = \{(G,s,\Gamma,h)\,\vert\, \hat{h} \leq h\text{ en }(G, \Gamma) \} \\
& \textbf{I}_{CON} = \{(G,s,\Gamma,h)\,\vert\, \hat{h}\text{ es consistente en }(G, \Gamma) \}
\end{split}
\end{equation}

Se puede demostrar que una heurística consistente es a su vez admisible, por lo que $\textbf{I}_{CON} \subseteq \textbf{I}_{AD}$. Además, de cada conjunto se toma un subconjunto especial denominado \textit{instancias no patológicas} que aquí no tomaremos en cuenta. Al conjunto de algoritmos que encuentran la solución óptima cuando hacen uso de una heurística admisible lo denotaremos como \textbf{A}$_{ad}$. De esta clase de algoritmos se estudian dos subclases: \textbf{A}$_{gc}$ representa a los algoritmos globalmente compatibles, esto es, aquellos que encuentran una solución óptima cuando A* lo hace, incluso cuando $\hat{h} > h$. \textbf{A}$_{bf}$ representa a los algoritmos de búsqueda primero el mejor admisibles.

\begin{figure}[H]
    \centering
    \begin{subfigure}[h]{.45\textwidth}
        \centering
        \includesvg[inkscapelatex=false, scale=0.35]{images/venn-instances.svg}
        \caption{}
    \end{subfigure}
    \hfill
    \begin{subfigure}[h]{.45\textwidth}
        \centering
        \includesvg[inkscapelatex=false, scale=0.35]{images/venn-classes.svg}
        \caption{}
    \end{subfigure}
    \caption{Diagramas de Venn de las instancias de problemas (a) y las clases de algoritmos considerados (b) \cite{dechter1985generalized}.}
    \label{fig:venn}
\end{figure}

En la Figura \ref{fig:venn} se pueden observar los diagramas de Venn correspondientes a los tipos de instancias de problemas y a las clases de algoritmos. Otro concepto necesario para definir la optimalidad es el de dominancia \cite{nilsson2014principles}:

\begin{defn}
Decimos que un algoritmo A domina a otro algoritmo B en un conjunto de instancias de problemas $\textbf{I}$ si y solo si, para cada instancia $I\in\textbf{I}$, el conjunto de nodos expandido por A es un subconjunto de los nodos expandidos por B. A domina estrictamente a B si y solo si A domina a B y B no domina a A. Esto es, existe al menos una instancia  donde A no expande un nodo que B sí, y no existe ninguna instancia donde ocurra lo contrario.
\end{defn}

\noindent Según este concepto de dominancia, se define la optimalidad de un algoritmo:

\begin{defn}
Un algoritmo A es óptimo sobre una clase $\textbf{A}$ de algoritmos si y solo si A domina a cada miembro de $\textbf{A}$.
\end{defn}

En base a esta definición, podemos clasificar la optimalidad de un algoritmo sobre otra clase de algoritmos con respecto a un tipo de instancias de problemas, pero debemos recordar que A* no es un único algoritmo, sino una familia de algoritmos diferenciados por lo que se conoce como regla de desempate. Esta regla de desempate es una forma de elegir el próximo nodo a expandir si en la lista abierta hay más de un nodo con el mismo valor de $f$. Teniendo esto en cuenta, las cuatro clases de optimalidad definidas para A* son las siguientes:

\begin{itemize}
    \item \textit{Tipo 0.} A* es 0-óptimo sobre \textbf{A} en \textbf{I} si y solo si, en cada instancia de problema, $I \in \textbf{I}$, cualquier regla de desempate expande un subconjunto de los nodos expandidos por cualquier miembro de \textbf{A}.
    \item \textit{Tipo 1.} A* es 1-óptimo sobre \textbf{A} en \textbf{I} si y solo si, para cada instancia de problema, $I\in\textbf{I}$, existe al menos una regla de desempate que expande un subconjunto de los nodos expandidos por cualquier miembro de $A$.
    \item \textit{Tipo 2.} A* es 2-óptimo sobre \textbf{A} en \textbf{I} si y solo si no existe ninguna instancia de problema $I\in\textbf{I}$ para la que algún miembro de \textbf{A} expande un subconjunto propio de los nodos expandidos por cualquier regla de desempate en A*.
    \item \textit{Tipo 3.} A* es 3-óptimo sobre \textbf{A} en \textbf{I} si y solo si se cumple que:
    Si existe una instancia de problema $I_1\in\textbf{I}$ donde algún algoritmo $B\in\textbf{A}$ no expande algún nodo expandido por alguna regla de desempate en A*, entonces existe alguna instancia $I_2\in\textbf{I}$ donde esa regla de desempate no expande un nodo expandido por B.
\end{itemize}

\begingroup
\setlength{\tabcolsep}{12pt} % Default value: 6pt
\renewcommand{\arraystretch}{2} % Default value: 1
\begin{table}[H]
    \centering
    \begin{tabular}{|c||c|c|c|}
    \hline
    & $\textbf{A}_{ad}$ & $\textbf{A}_{gc}$ & $\textbf{A}_{bf}$ \\
    \hline\hline
    $\textbf{I}_{AD}$ & No óptimo & 1-óptimo & 1-óptimo \\
    \hline
    $\textbf{I}_{CON}$ & 1-óptimo & 1-óptimo & 1-óptimo \\
    \hline
    \end{tabular}
    \caption{Tipo de optimalidad de A* sobre otras clases de algoritmos en distintas instancias de problemas \cite{dechter1985generalized}.}
    \label{tab:optimality}
\end{table}
\endgroup

En la Tabla \ref{tab:optimality} podemos observar el tipo de optimalidad de A* en comparación con las tres clases de algoritmos en los cuatro dominios de problemas. Los resultados más destacables se concentran en la primera columna: A* no es óptimo sobre $\textbf{A}_{ad}$ para aquellos problemas en los que la heurística solo es admisible, es más, ningún algoritmo óptimo puede existir a menos que la heurística sea también consistente. Recordemos por los tipos de optimalidad que esto quiere decir que pueden existir algoritmos que consiga superar a A* en algunas instancias de esta clase de problemas. Por la segunda y tercera columna podemos ver que A* sí que es óptimo sobre las clases de algoritmos $\textbf{A}_{gc}$ y $\textbf{A}_{bf}$.

Ahora que ya hemos comprendido el funcionamiento del algoritmo A*, su condición de admisibilidad y su optimalidad, podemos abordar las distintas estrategias que se han propuesto para paralelizarlo. En la Sección \ref{sec:estado_arte} veremos las distintas estrategias de paralelización y cómo se relacionan con los modelos de computación paralela. Comenzaremos describiendo las estrategias centralizadas y posteriormente las descentralizadas, haciendo especial hincapié en las estrategias basadas en distribución de trabajo con funciones \textit{hash}. Acto seguido, en la Sección \ref{sec:objetivos} se detallarán los objetivos propuestos para el resto del trabajo y en la Sección \ref{sec:metodologia} revisaremos la metodología seguida para conseguirlos. Se explicarán tres implementaciones del algoritmo y cómo se han conseguido para, finalmente, comparar su rendimiento en distintos tipos de escenarios en las Secciones \ref{sec:implementacion} y \ref{sec:resultados}. Por último, en la Sección \ref{sec:conclusiones} se presentarán las conclusiones y en la Sección \ref{sec:futuro} se propondrán varias líneas de trabajo futuro que pueden surgir a partir de este estudio.

\newpage
%%%%%%%%%%%%%%%%%%%%%%%%%%%%%%%%%%%%%%%%%%%%%%%%%%%%%%%%%%%%%%%%%%%%%%%%%%%%%%%
\section{Estado del arte} \label{sec:estado_arte}

El algoritmo A* es eficaz para encontrar caminos óptimos en grafos, pero su rendimiento se ve limitado cuando se aplica a problemas de gran escala, debido a la gran cantidad de memoria que requiere para mantener la información de los nodos visitados. Para mitigar este problema, además de las modificaciones limitadas en memoria de la versión secuencial, se han propuesto diversas estrategias para paralelizar el algoritmo.

Las primeras propuestas para paralelizar un algoritmo de búsqueda aparecen en las décadas de los 70 y 80. Wyllie \cite{wyllie1979complexity} propone en su tesis doctoral paralelizar la estrategia de búsqueda en anchura y, aunque otros autores desestimaron su idea afirmando que no era posible \cite{eckstein1977parallel, reghbati1978parallel, reif1985depth}, a principios de la década de los 80 comenzaron a surgir avances en éste área y que podemos encontrar en orden cronológico en el artículo de Freeman \cite{freeman1991parallel}, demostrando la viabilidad de la idea. Aunque estas ideas no incorporaban la información heurística en el problema como A*, sentaron las bases para muchos de los principios que se aplicarían después. Rao y Kumar \cite{rao1987parallel, kumar1987parallel} también trabajaron con la búsqueda en profundidad y, como resultado de su trabajo, unieron la búsqueda en profundidad con el algoritmo IDA*, consiguiendo una aumento en la velocidad superlineal \cite{rao1988superlinear}. 

Estos primeros trabajos se enfocaban en dividir el espacio de búsqueda entre distintos procesos\footnote{En la literatura es habitual referirse a proceso cuando se habla la paralelización del algoritmo. En este trabajo hemos mantenido esta nomenclatura, aunque cabe destacar que cuando se implementa en un entorno de memoria compartida como OpenMP, lo correcto es hablar de hilo, y así haremos en la sección dedicada al diseño y la implementación de nuestra solución.}, cada uno de los cuales exploraba una región diferente del grafo o árbol de búsqueda. En general, se trataba de problemas en los que el espacio de estados se puede representar como una estructura tipo árbol. De ésta manera, los subárboles generados a partir de cada nodo son independientes, y la paralelización es casi inmediata. Sin embargo, en la mayoría de aplicaciones, los subárboles generados son irregulares, provocando un desbalanceo en la carga de trabajo entre los procesos.

% Antes de abordar las técnicas concretas de paralelización de A*, es necesario introducir brevemente los principales modelos de computación paralela que, además, representan fielmente las dos principales estrategias con las que se intenta paraleliza el algoritmo. En primer lugar, tenemos el modelo de memoria compartida, en el que varios procesos se ejecutan en paralelo y acceden a una memoria común, lo que facilita el intercambio de información pero requiere mecanismos de sincronización para evitar condiciones de carrera. Por otro lado, tenemos el modelo de memoria distribuida, que se basa en procesos que no comparten memoria y se comunican entre sí mediante el envío de mensajes. Este segundo modelo implica un diseño más explícito de la coordinación entre procesos, pero ofrece una mejor escalabilidad en arquitecturas con muchos núcleos o nodos.

Al diseñar algoritmos de búsqueda paralelos, es importante tener en cuenta tres problemas principales: la sobrecarga de búsqueda, que se refiere a las implementaciones paralelas que expande más nodos que una implementación secuecnial; la sobrecarga de sincronización, que se refiere a cuando un proceso está a que otros lleguen a un punto de sincronización; y la sobrecarga de comunicación, que tiene que ver con el coste de compartir información entre los procesos. Las diferentes estrategias que se han propuesto en la literatura se diferencian principalmente en cómo tratan estos problemas. Kumar \cite{kumar1988parallel} propone dos enfoques para paralelizar las búsquedas primero el mejor, diferenciados por cómo usan y gestionan la lista abierta entre los diferentes procesos. En la revisión de Fukunaga \cite{fukunaga2017survey} encontramos una imagen que se ha traducido al castellano para este trabajo y que representa muy bien las distintas estrategias y a qué algoritmos han dado lugar (Figura \ref{fig:clasficiacion}).

\begin{figure}[H]
    \centering
    \includesvg[inkscapelatex=false, scale=0.7]{images/classification.svg}
    \caption{Clasificación de las distintas estrategias de paralelización \cite{fukunaga2017survey}.}
    \label{fig:clasficiacion}
\end{figure}


\subsection{A* centralizado} \label{sec:centralizado}

La idea más sencilla para paralelizar A* es extraer tantos nodos de la lista abierta como procesos estén participando en la búsqueda. Esta estrategia es conocida como la estrategia centralizada, pues tanto la lista abierta como la cerrada se mantiene en una memoria global compartida por todos los procesos (Figura \ref{fig:centralizado}). Este enfoque presenta dos desafíos principalmente, el primero es que el criterio de terminación del algoritmo secuencial deja de ser válido. No podemos asegurar que cuando un nodo es extraído de a lista abierta, éste sea la mejor solución, pues puede haber otro proceso que en su próxima expansión encuentra una solución mejor. Esto ocurre cuando hay más de un nodo objetivo, y el criterio de terminación puede modificarse para que el algoritmo termine únicamente cuando encuentra la mejor solución \cite{leifker1984hybrid}. Por otro lado, los continuos accesos a la lista abierta generan una zona de contención entre los procesos, ya que ésta debe protegerse mediante un cerrojo para evitar condiciones de carrera. Esta contención limitaría el \textit{speedup} obtenido en la versión paralela a $(T_{exp} + T_{access})/T_{access}$, donde $T_{exp}$ es el tiempo promedio para una expansión, y $T_{access}$ es el tiempo promedio accediendo a la lista abierta por cada expansión \cite{huang1987tight}. También es necesario proteger la lista cerrada, pues diferentes procesos pueden estar leyendo y escribiendo a la vez para comprobar si el nodo ya ha sido visitado o si es un nodo nuevo. Lo más habitual es establecer un cerrojo por nodo, de manera que no se bloquee la estructura entera para el resto de procesos que puedan estar expandiendo otros nodos diferentes.

\begin{figure}[H]
    \centering
    \includesvg[inkscapelatex=false, width=\linewidth]{images/centralized.svg}
    \caption{Estrategia centralizada con lista abierta compartida en memoria global. Cada proceso bloque la lista, inserta los nodos expandidos en la iteración anterior, extrae el de menor coste y la desbloquea para el siguiente proceso \cite{grama1995parallel}.}
    \label{fig:centralizado}
\end{figure}

Este algoritmo recibe el nombre de \textit{Simple Parallel A*} o SPA* (Algoritmo \ref{spastar}). En el pseudocódigo no se muestran los cerrojos de la lista cerrada y una variable atómica extra que es necesaria para detectar la terminación. Se comentarán estas diferencias a lo largo de la explicación, pero se ha decidido no incluirlas para no hacerlo excesivamente largo y difícil de entender a simple vista.

La estructura general es similar a A*, inicializa la lista abierta con el nodo inicial y la lista cerrada vacía, pero ahora cuenta con un cerrojo para la lista abierta y otro para $m$, que representa la mejor solución encontrada hasta ahora (líneas 1-4). En esta parte también se inicializarían los cerrojos de la lista cerrada y el contador atómico. El bucle principal del algoritmo se ejecuta en paralelo para cada proceso $p$ (líneas 6-41). En esta versión no podemos terminar el algoritmo cuando $Q\neq \emptyset$ (línea 7), ya que puede darse el caso de que la lista esté vacía y otro proceso esté expandiendo un nodo sin haber introducido todavía los vecinos en la lista, algo que suele ocurrir muy a menudo al comienzo del algoritmo. Además, se debe comprobar también que el próximo nodo no empeora el coste del objetivo actual, pues se harían expansiones innecesarias. El valor estimado para $m$ cuando aún no hemos encontrado el objetivo debería ser un valor lo suficientemente alto, esto es, $\hat{f}(\text{NULL})=\infty$. Cuando se cumple alguna de las condiciones, se debe comprobar que ningún proceso puede descubrir una mejor solución.

\begin{algorithm}[H]
\begin{algorithmic}[1]
\small
\caption{Simple Parallel A* (SPA*)}\label{spastar}
\State $Q \gets \{s\}$;
\State $H \gets \emptyset$;
\State Lock $l_o, l_m$;
\State $m \gets$ NULL;
\State En paralelo, para cada proceso $p$, ejecutar:
\While{$\textsc{TerminateDetection}()$}
    \If{$Q = \emptyset$ \textbf{or} $\textsc{PeekMin}(Q) \geq \hat{f}(m)$}
        \State \textbf{continue};
    \EndIf
    \State $\textsc{AcquireLock}(l_o)$;
    \State $n \gets \textsc{ExtractMin}(Q)$;
    \State $\textsc{ReleaseLock}(l_o)$;
    \State $\textsc{Insert}(H, n)$;
    \If{$n = t$}
        \State $\textsc{AcquireLock}(l_m)$;
        \If{$\hat{g}(n) < \hat{g}(m)$}
            \State $m \gets n$;
        \EndIf
        \State $\textsc{ReleaseLock}(l_m)$;
    \EndIf
    \ForAll{$n' \in \textsc{Neighbors}(n)$}
        \State $g' \gets \hat{g}(n)+c(n, n')$;
        \If{$n' \in H$}
            \If{$g' < \hat{g}(n')$}
                \State $\textsc{Remove}(H,n')$;
                \State $\textsc{Insert}(Q, n')$;
            \Else
                \State \textbf{continue};
            \EndIf
        \Else
            \If{$n' \in Q$}
                \State $\textsc{Insert}(Q,n')$;
            \ElsIf{$g_1 \geq \hat{g}(n')$}
                \State \textbf{continue};
            \EndIf
        \EndIf
        \State $\hat{g}(n') \gets g'$;
        \State $\hat{f}(n') \gets \hat{g}(n') + \hat{h}(n')$;
        \State $parent(n') \gets n$;
    \EndFor
\EndWhile
\If{$m =$ NULL}
    \State \textbf{return} error
\Else
    \State \textbf{return} \textsc{RetracePath}($m$);
\EndIf
\end{algorithmic}
\end{algorithm}

Para detectar la correcta terminación del algoritmo se utiliza variable atómica que cada proceso aumenta cuando se cumple alguna de las condiciones y que vuelve a decrementar si no se cumplen en alguna iteración posterior. Al inicio de cada iteración del bucle, \textsc{TerminateDetection()} se encarga de comprobar si el valor de esta variable es igual al número de procesos, en cuyo caso se puede dar por terminado el algoritmo. 

A continuación los procesos se disponen a extraer el siguiente nodo de la lista abierta, protegiendo el acceso con un cerrojo para evitar condiciones de carrera (líneas 10-12). La inserción en la lista cerrada (línea 13) puede hacerse sin ningún cerrojo, ya que cada proceso habrá extraído un nodo diferente. La comprobación del objetivo es similar a la de la versión secuencial, pero protegiendo la variable $m$ con un cerrojo (líneas 14-20). El resto del algoritmo, es decir, la expansión del nodo actual, se realiza igual que en el algoritmo secuencial (líneas 21-41), pero en este caso, tanto los accesos a $Q$ como a $H$ deben protegerse. Puede darse el caso de que dos procesos obtengan el mismo vecino desde dos nodos diferentes, siendo necesaria la sincronización a la hora de comprobar cuál de los dos debe actualizar los valores del nodo si éste está cerrado y se han encontrado mejores caminos. Cuando el bucle principal termina, se comprueba si se ha encontrado el nodo objetivo y, si es el caso, se devuelve el camino (líneas 42-46).

SPA* es muy sencillo de implementar en un modelo de memoria compartida, pero ha probado ser poco eficiente \cite{lai1984anomalies}, incluso peor que la versión secuencial para problemas en los que la operación de expansión de un nodo es rápida, debido a que la mayor parte del tiempo de ejecución de los procesos se desperdicia mientras esperan a obtener un nodo de la lista abierta. Aun así, ha servido de base para otras propuestas como PA*SE \cite{phillips2014pa}, que introduce un mecanismo para evitar reexpansiones innecesarias en SPA*, o PKBFS \cite{vidal2010adaptive}, que realiza la expansión paralela de los $k$ mejores nodos en cada iteración.

\subsection{A* descentralizado} \label{sec:descentralizado}

La estrategia descentralizada surge para aliviar la sobrecarga debida a los accesos constantes a las listas abierta y cerrada en SPA*. En contraste con la estrategia centralizada, cada proceso tiene su propia lista abierta. Al inicio, el proceso padre expande el primer nodo y reparte los sucesores entre el resto de procesos hijos, que empezará una búsqueda en cada uno de ellos. Al descentralizar la lista abierta, se eliminan las condiciones de carrera al intentar acceder a una estructura global, pero surge la necesidad de balancear la carga de trabajo.

En una estrategia distribuida, cada proceso extrae un nodo de su lista local. Cuando el espacio de búsqueda tiene forma de árbol, esta estrategia divide el trabajo de forma equitativa entre los procesos, pero puede llevar a que alguno de ellos explore una parte del árbol que hubiese sido descartada por la versión secuencial. Varios autores han propuesto estrategias de comunicación para repartir los nodos entre los procesos y aliviar esta sobrecarga en la búsqueda, lo que representa el clásico intercambio entre computación y comunicación en las arquitecturas distribuidas. Cuando el espacio de búsqueda es un grafo, además, surge el problema de detectar la duplicación de nodos, es decir, dos procesos pueden obtener el mismo nodo como resultado de la expansión de dos nodos diferentes.

% {\color{blue} revisa errores tipograficos y gramaticale en textos anteriores} % 


En una estrategia descentralizada (Algoritmo \ref{decentralized-astar}), cada proceso cuenta con su propia lista abierta (línea 1), cerrada y un \textit{buffer} de recepción. La función \textsc{ComputeRecipient}$(n)$ devuelve el índice del proceso al que debe mandarse el nodo $n$. Al entrar en el bucle principal (líneas 5-42), cada proceso comprueba si hay algún elemento en el \textit{buffer}. Extrae la información sobre el nodo, su nuevo coste y su padre (línea 7). A continuación comprueba si es un duplicado, si mejora el coste anterior, si es un nuevo nodo o si se puede descartar, al igual que en las otras versiones centralizada y secuencial (líneas 8-24). Una vez vaciado el \textit{buffer}, comprueba si debe terminar (líneas 26-28). Extrae un nodo de su lista abierta y lo inserta en su lista cerrada (líneas 29-30). Si es un nodo objetivo, actualiza el valor de $m$ (líneas 31-37). En caso contrario, expande el nodo, calculando el nuevo coste para alcanzarlo y enviando al propietario la información.

Gran parte de la eficiencia de estas estrategias depende de la función \textsc{ComputeRecipient(n)} y de cómo reparta los nodos entre procesos. Karp y Zhang \cite{karp1988randomized, karp1993randomized} proponen una estrategia de distribución de trabajo aleatoria en la que cada vez que se obtiene un nodo al expandir, se calcula a qué proceso debe asignarse. Kumar et al. \cite{kumar1988parallel} proponen otro tipo de estrategia de distribución aleatoria a la que denominan \textit{blackboard}, en la que los procesos comparten información a través de un estructura común. Otro ejemplo de estrategia distribuida es la denominada búsqueda bidireccional, que consiste en la ejecución del algoritmo de búsqueda desde el nodo inicial hacia el final y desde el nodo final hacia el inicial \cite{rios2011pnba}.

El problema con las estrategias aleatorias es que los nodos duplicados solo son detectados si se mandan fortuitamente al mismo proceso. Otro enfoque más simple para distribuir la carga de trabajo entre los procesos fue introducido por Evett et al. \cite{evett1995massively}. En su artículo presentan PRA*, una versión paralela de A* distribuida en la que el reparto de nodos entre los procesos se realiza mediante una función \textit{hash} que asigna cada nodo a único proceso. Al expandir un nodo, cada proceso debe calcular a qué proceso debe mandar cada sucesor de manera síncrona. Cada proceso cuenta con un \textit{buffer} que el resto deben bloquear cada vez que envían un nodo. Esta idea de asignar el trabajo a los procesos mediante una función \textit{hash} también fue usada por Mahapatra \cite{mahapatra2002scalable} para paralelizar una versión de A* llamada SEQ\_A*, pero tanto Evett et al. como Mahapatra no explorar esta técnica lo suficiente en sus trabajos, ya que se centraban en otros aspectos del problema, como la retracción de nodos en el caso de Evett.

\begin{algorithm}[H]
\caption{A* descentralizado}\label{decentralized-astar}
\begin{algorithmic}[1]
\small
\State $Q \gets \{Q_i\}_{i=1}^p$;
\State $\textsc{Insert}(Q_{\textsc{ComputeRecipient}(s)}, s)$;
\State $m \gets$ NULL;
\State En paralelo, para cada proceso $p$, ejecutar:
\While{$\textsc{TerminateDetection}()$}
    \While{$BUFFER_P \neq \emptyset$}
        \State $(n', g', n) \gets \textsc{Extract}(BUFFER_p)$;
        \If{$n'\in H_p$}
            \If{$g' < \hat{g}(n')$}
                \State $\textsc{Remove}(H_p, n')$;
                \State $\textsc{Insert}(Q_p, n')$;
            \Else
                \State \textbf{continue};
            \EndIf
        \Else
            \If{$n' \notin Q_p$}
                \State $\textsc{Insert}(Q_p, n')$;
            \ElsIf{$g' \geq \hat{g}(n')$}
                \State \textbf{continue};
            \EndIf
        \EndIf
        \State $\hat{g}(n') \gets g'$;;
        \State $\hat{f}(n') \gets \hat{g}(n') + \hat{h}(n')$;
        \State $parent(n') \gets n$;
    \EndWhile
    \If{$Q_p = \emptyset$ \textbf{or} $\textsc{PeekMin}(Q_p) \geq \hat{f}(m)$}
        \State \textbf{continue};
    \EndIf
    \State $n \gets \textsc{ExtractMin}(Q_p)$;
    \State $\textsc{Insert}(H_p, n)$;
    \If{$n = t$}
        \State $AcquireLock(l_m)$
        \If{$\hat{g}(n) < \hat{g}(m)$}
            \State $m \gets n$;
        \EndIf
        \State $ReleaseLock(l_m)$;
    \EndIf
    \ForAll{$n' \in \textsc{Neighbors}(n)$}
        \State $g' \gets \hat{g}(n) + c(n, n')$;
        \State $\textsc{Insert}(BUFFER_{\textsc{ComputeRecipient}(n')}, (n', g', n))$;
    \EndFor
\EndWhile
\If{$m =$ NULL}
    \State \textbf{return} \textsc{error}
\EndIf
\State \textbf{return} \textsc{RetracePath}($m$);
\end{algorithmic}
\end{algorithm}

\subsection{Hash Distributed A*}

Kishimoto, Fukunawa y Botea \cite{kishimoto2009scalable, kishimoto2013evaluation} retoman esta idea de repartir los nodos entre procesos mediante una función \textit{hash} años más tarde y proponen HDA* (Hash-Distributed A*), una versión simple, paralela y escalable de A*, en la que a cada proceso le pertenece una parte del espacio de búsqueda. HDA* opera de manera similar a la expuesta en el Algoritmo \ref{decentralized-astar}. La principal diferencia con la propuesta de PRA* es que la comunicación entre los procesos se realiza de manera asíncrona, lo que evita tener que bloquear los \textit{buffers} de recepción del resto de procesos cuando se le deben enviar nuevos nodos. Las primeras implementaciones de PRA* mostraban un rendimiento muy bajo, incluso peor que la versión secuencial, en problemas como la búsqueda en \textit{grids} \cite{burns2010best} debido a la sobrecarga introducida por los bloqueos. En su lugar, en HDA*, la comunicación entre procesos se implementa mediante paso de mensajes incluso en entornos de memoria compartida.

Los primeros trabajos sobre HDA* no evaluaron de forma cuantitativa el impacto que tiene la elección de la función \textit{hash}, lo que llevó a que investigaciones posteriores arrojaran resultados poco precisos. Posteriormente, Jinnai y Fukunaga \cite{jinnai2016abstract} realizaron una comparación entre diferentes funciones de dispersión utilizadas en la literatura, demostrando que tanto la función \textit{hash} de Zobrist como su versión mejorada, conocida como Abstract Zobrist hashing, superan con claridad a otras alternativas empleadas anteriormente.


\subsection{Implementaciones en GPU}

En los últimos años se han desarrollado implementaciones de A* en arquitecturas de procesamiento gráfico (GPU). Zheng y Zhou \cite{zhou2015massively} proponen GA*, un algoritmo que explota el paralelismo para la expansión de nodos utilizando grandes cantidades de listas abiertas distribuidas entre los hilos. Esto provoca la generación de una gran cantidad de nodos duplicados, pero proponen dos técnicas efectivas basadas en \textit{hashing} para detectarlos. La principal limitación de esta implementación es la limitada memoria de las GPU, lo que ha llevado a otras propuestas basadas en versiones limitadas en memoria como IDA* \cite{horie2017block}.

\newpage
%%%%%%%%%%%%%%%%%%%%%%%%%%%%%%%%%%%%%%%%%%%%%%%%%%%%%%%%%%%%%%%%%%%%%%%%%%%%%%%
\section{Objetivos} \label{sec:objetivos}
%%%%%%%%%%%%%%%%%%%%%%%%%%%%%%%%%%%%%%%%%%%%%%%%%%%%%%%%%%%%%%%%%%%%%%%%%%%%%%%

El objetivo principal de este trabajo es analizar el impacto del paralelismo en el rendimiento del algoritmo A*. Para ello, se implementará una versión secuencial y dos paralelas: una centralizada y otra descentralizada, ambas implementadas en un modelo de memoria compartida. A través de esta comparación, se pretende conseguir una mejora en el tiempo de ejecución con respecto a la versión secuencial, así como estudiar las limitaciones y desafíos asociados a ambos enfoques de paralelización.

En concreto, se plantean los siguientes objetivos específicos:

\begin{itemize}
\item Implementar una versión secuencial del algoritmo A* que sirva como base para las comparaciones.

\item Desarrollar una versión paralela centralizada del algoritmo basada en SPA*.

\item Implementar una versión paralela descentralizada basada en el algoritmo HDA* y adaptada a un entorno de memoria compartida.

\item Evaluar el rendimiento de cada versión en términos de tiempo de ejecución, escalabilidad y eficiencia en distintos tipos grafos.

\item Analizar los resultados obtenidos y extraer conclusiones sobre la viabilidad de los dos enfoques principales de paralelización del algoritmo.
\end{itemize}

\newpage

%%%%%%%%%%%%%%%%%%%%%%%%%%%%%%%%%%%%%%%%%%%%%%%%%%%%%%%%%%%%%%%%%%%%%%%%%%%%%%%
\section{Metodología} \label{sec:metodologia}
%%%%%%%%%%%%%%%%%%%%%%%%%%%%%%%%%%%%%%%%%%%%%%%%%%%%%%%%%%%%%%%%%%%%%%%%%%%%%%%

Para abordar los objetivos planteados, se han implementado las disitintas versiones usando el lenguaje de programación C y el compilador GCC. Para la implementación del paralelismo y la gestión de la memoria compartida se ha utilizado la librería OpenMP.

La primera fase del trabajo consistió en la implementación de una versión secuencial, la cual actúa como línea base para evaluar los beneficios de las versiones paralelas.

La versión paralela centralizada se basa en el enfoque SPA* (Simple Parallel A*), donde múltiples hilos de OpenMP comparten las listas abierta y cerrada y cooperan en la expansión de nodos. La sincronización y la gestión de la concurrencia se realizan mediante directivas de OpenMP, lo que permite aprovechar el paralelismo sin salir del entorno de memoria compartida.

La versión descentralizada se inspira en el algoritmo HDA* (Hash Distributed A*), originalmente diseñado para entornos distribuidos y programado típicamente con MPI. En este trabajo, se ha adaptado dicho enfoque para funcionar en un entorno compartido, simulando la comunicación asíncrona entre procesos mediante mecanismos de envío dinámico de nodos. Esta adaptación permite estudiar las características del paralelismo descentralizado sin la necesidad de implementar una infraestructura distribuida real.

El rendimiento de las distintas versiones se ha realizado con las utilidades que provee la propia librería OpenMP para extraer información del tiempo en ejecución. Se ha seleccionado un conjunto de escenarios de prueba que incluyen grafos de diferentes tamaños y densidades. 

Los resultados obtenidos nos muestran las ventajas y limitaciones de cada enfoque de paralelización. Se estudian el tiempo de ejecución de cada versión al variar el número de hilos y el tipo de grafo con el fin de analizar la escalabilidad y eficiencia de cada enfoque.

\newpage

%%%%%%%%%%%%%%%%%%%%%%%%%%%%%%%%%%%%%%%%%%%%%%%%%%%%%%%%%%%%%%%%%%%%%%%%%%%%%%%
\section{Diseño e implementación} \label{sec:implementacion}
%%%%%%%%%%%%%%%%%%%%%%%%%%%%%%%%%%%%%%%%%%%%%%%%%%%%%%%%%%%%%%%%%%%%%%%%%%%%%%%

En esta sección se explicará en detalle el proceso seguido para obtener tanto la versión secuencial como las versiones paralelas a partir de ésta. Finalmente, compararemos los tres tipos de implementaciones en distintos escenarios y analizaremos los tiempos obtenidos.

La implementación cuenta con los siguientes ficheros:

\begin{itemize}
    \item \verb|astar.h|: Cabecera con las estructuras y métodos comunes a las tres implementaciones. También define el método para iniciar la versión secuencial.
    \item \verb|astar.c|: Implementación secuencial.
    \item \verb|spastar.h|: Cabecera que define el método para iniciar SPA*.
    \item \verb|spastar.c|: Implementación de SPA*.
    \item \verb|hdastar.h|: Cabecera que define las estructuras necesarias para la versión descentralizada. También define el método para iniciar HDA*.
    \item \verb|hdastar.c|: Implementación de HDA*.
\end{itemize}

Cabe mencionar antes de comenzar con las implementaciones una modificación que es muy frecuente encontrar en las implementaciones más modernas de A*. Se trata de utilizar una lista de visitados en lugar de una lista cerrada. La diferencia fundamental es que cuando un nodo es reabierto, no se elimina de la lista de visitados. Esta situación, que se da cuando la heurística es admisible pero no consistente, conlleva en la versión original eliminar al nodo de la lista cerrada e insertarlo de nuevo en la lista abierta. El problema que esto presenta es que si el nodo no está cerrado pero sí está abierto, la operación de buscarlo en la lista abierta es bastante costosa: $O(n)$ si se implementa como un montículo binario. La solución pasa por no eliminar el nodo de la lista cerrada y añadir un campo a la estructura que representa a los nodos que nos indique si el nodo está abierto o no. Al no eliminar nunca los nodos, esta pasa a llamarse lista de visitados.

Esta modificación conlleva un cambio en el pseudocódigo del Algoritmo \ref{astar}. En concreto, no se eliminaría el nodo de la lista cerrada en la línea 13 y las líneas 18 a 23 serían sustituidas por la inserción del nodo en la lista de visitados y en la lista de abiertos. Se puede ver esta modificación en el Algoritmo \ref{visited-astar}, donde $V$ representa la lista de visitados.

\begin{algorithm}[H]
\caption{Modificación de A* (Algoritmo \ref{astar}, líneas 9-28)}\label{visited-astar}
\begin{algorithmic}[1]
\setcounter{ALG@line}{8}
\ForAll{$n' \in \textsc{Neighbors}(n)$}
    \State $g' \gets \hat{g}(n)+c(n, n')$;
    \If{$n' \in V$}
        \If{$g' < \hat{g}(n')$}
            \State $\textsc{Insert}(Q, n')$;
        \Else
            \State \textbf{continue};
        \EndIf
    \Else
        \State $\textsc{Insert}(V, n')$
        \State $\textsc{Insert}(Q, n')$;
    \EndIf
    \State $\hat{g}(n') \gets g'$;
    \State $\hat{f}(n') \gets g(n') + h(n')$;
    \State $parent(n') \gets n$;
\EndFor
\end{algorithmic}
\end{algorithm}

\subsection{Implementación secuencial}

En primer lugar veremos la implementación de la versión secuencial. Empezaremos por las estructuras de datos utilizadas. Para representar a los nodos hemos utilizado una estructura \verb|node_t| que podemos encontrar en \verb|astar.h|:

\begin{lstlisting}
struct {
    int id;
    float gCost;
    float fCost;
    int parent;
    int open_index;
} node_t;
\end{lstlisting}

Cuenta con los valores de $\hat{g}(n)$, $\hat{f}(n)$ y $parent(n)$ que hemos visto antes y, además, con un identificador que nos permite diferenciar a los nodos y un campo \verb|open_index| que nos indica la posición en la lista abierta. Este valor será $-1$ cuando el nodo no esté abierto. Para iniciar los algoritmos es necesario definir una estructura \verb|AstarSource|. Esta estructura cuenta con tres parámetros: el número máximo de nodos en el espacio de búsqueda, \verb|max_size|, y dos punteros a funciones. La primera de ellas recibe como parámetros los identificadores de dos nodos y devuelve la estimación de la distancia entre ellos, la heurística. La segunda función se encarga de calcular los vecinos. Ambas deben ser implementadas por el usuario para el dominio específico en el que se quiera ejecutar los algoritmos.

\begin{lstlisting}
struct {
    int max_size;
    float (*heuristic)(int id_1, int id_2);
    void (*get_neighbors)(neighbors_list *neighbors, int id);
} AStarSource;
\end{lstlisting}

Al definir la función \verb|get_neighbors|, el usuario debe añadir los nodos vecinos con la función:

\begin{lstlisting}
void add_neighbor(neighbors_list *neighbors, int n_id, float cost);
\end{lstlisting}

\noindent indicando el índice del vecino, \verb|n_id|, y el coste para alcanzarlo,\verb|cost|. La lista de vecinos, \verb|neighbors|, será la que recibe como parámetro en \verb|get_neighbors|. El algoritmo irá haciendo llamadas a este método en cada iteración para obtener la lista de vecinos del nodo actual. La estructura de la lista de vecinos es la siguiente:

\begin{lstlisting}
struct {
    int capacity;
    int count;
    int *nodeIds;
    float *costs;
} neighbors_list;
\end{lstlisting}

Para iniciar la búsqueda búsqueda con el algoritmo secuencial, se debe hacer una llamada al siguiente método, que recibe como parámetro la estructura \verb|AStarSource| y los identificadores del nodo inicial y el final:

\begin{lstlisting}
path *astar_search(AStarSource *source, int s_id, int t_id);
\end{lstlisting}

\noindent que devuelve un puntero a una estructura \verb|path| en caso de encontrar el camino, o \verb|NULL| si no existe el camino. La estructura que representa los caminos tiene la siguiente forma:

\begin{lstlisting}
struct {
    int count;
    int *nodeIds;
    float cost;
} path;
\end{lstlisting}

\noindent donde \verb|count| representa el número de nodos en el camino, \verb|nodeIds| es la lista de identificadores de dichos nodos, en orden, y \verb|cost| es el coste total del camino.

En cuanto a las estructuras de datos propias del algoritmo, la lista abierta se ha implementado con un montículo binario en los ficheros \verb|heap.h| y \verb|heap.c|. Las tres versiones utilizan la misma implementación. La lista cerrada consiste en un array de punteros a estructuras \verb|node_t| y se inicializada con el tamaño especificado por \verb|max_size|. Ambas estructuras, junto con la lista de vecinos, se crean al comienzo del algoritmo.

\subsection{Implementación centralizada}

La versión paralela centralizada sigue la misma estructura que la versión secuencial. En este caso se utilizan estructuras \verb|omp_lock_t| para proteger la lista abierta y cada uno de los nodos de la lista de visitados.

\begin{lstlisting}
omp_lock_t open_lock;
omp_lock_t *visited_locks;
\end{lstlisting}

Cada vez que un hilo va a extraer un nodo de la lista abierta debe adquirir el cerrojo, y en el caso de encontrar un mejor camino para uno de los nodos ya en la lista de visitados (abierto o no) también debe bloquear ese nodo. OpenMP ofrece estos mecanismos mediante dos métodos que reciben como parámetro la dirección del cerrojo a adquirir o liberar:

\begin{lstlisting}
omp_set_lock(omp_lock_t *lock);
omp_unset_lock(omp_lock_t *lock);
\end{lstlisting}

La sección concurrente se declara justo antes del bucle \verb|while| principal, con la directiva OpenMP:

\begin{lstlisting}
# pragma omp parallel if(k > 1) num_threads(k) shared(open, open_lock, visited, visited_locks, m, m_lock, terminated, terminated_lock)
\end{lstlisting}

Las variables compartidas entre hilos corresponden, en el orden en que aparecen en la directiva, a la lista abierta y su cerrojo, la lista de visitados y sus cerrojos, la mejor solución actual y su cerrojo, el contador de terminación y su cerrojo. Esta sección solamente se ejecutará de manera concurrente si el número de hilos es mayor a uno, que se especifica como parámetro al iniciar la búsqueda:

\begin{lstlisting}
path *spastar_search(AStarSource *source, int start_id, int goal_id, int k)
\end{lstlisting}

Las principales diferencias de esta versión con la secuencial son dos. En primer lugar, la adquisición y liberación de cerrojos al extraer nodos de la lista abierta y en la evaluación de los vecinos de cada nodo, como ya hemos comentado. La segunda diferencia es la condición de terminación que podemos implementar mediante el uso de un contador atómico (Sección \ref{sec:centralizado}). Este contador es la variable \verb|terminated|, y los accesos atómicos se controlan mediante el cerrojo \verb|terminated_lock|, aunque se podrían haber encapsulado también las modificaciones dentro de una directiva \verb|critical| de OpenMP. El código para la detección de la terminación (Código \ref{terminacion}) se ejecuta al comenzar cada iteración del bucle principal. En primer lugar se comprueba si todos los hilos han terminado (líneas 2-7). Si no es así, se obtiene el valor del mejor camino encontrado actualmente (líneas 9-11) con exclusión mutua mediante el cerrojo \verb|m_lock|. Se debe comprobar si la lista abierta está vacía o si el próximo nodo a extraer es peor que la mejor solución encontrada hasta ahora (línea 14). Esta operación también se hace con exclusión mutua sobre la lista abierta, pues puede haber otros hilos insertando nodos en ese momento. Si se cumple alguna de las condiciones, el hilo aumenta el valor de la variable \verb|terminated| también con exclusión mutua. Se utiliza una variable auxiliar privada para cada hilo, \verb|waiting|, que evita que el mismo hilo aumente el contador varias veces (líneas 16-21). Finalmente, si no se cumple ninguna de las condiciones el hilo puede continuar, pero si había aumentado el contador anteriormente, ahora debe decrementarlo (líneas 25-30).

{\lstset{frame=trbl,
         numbers=left,
         caption=Detección de terminación en A* centralizado,
         label=terminacion}
\begin{lstlisting}
...
omp_set_lock(&terminated_lock);
if (terminated == k) {
    omp_unset_lock(&terminated_lock);
    break;
}
omp_unset_lock(&terminated_lock);

omp_set_lock(&m_lock);
float m_cost = (m != NULL) ? m->fCost : FLT_MAX;
omp_unset_lock(&m_lock);

omp_set_lock(&open_lock);
if (heap_is_empty(open) || heap_min(open) >= m_cost) {
    omp_unset_lock(&open_lock);
    if (!waiting) {
        waiting = 1;
        omp_set_lock(&terminated_lock);
        terminated++;
        omp_unset_lock(&terminated_lock);
    }
    continue;
}

if (waiting) {
    waiting = 0;
    omp_set_lock(&terminated_lock);
    terminated--;
    omp_unset_lock(&terminated_lock);
}
...
\end{lstlisting}
}

\clearpage

\subsection{Implementación descentralizada}

La versión descentralizada es la que más se diferencia de las otras dos en cuanto a la implementación se refiere. El bucle principal no es tan paracedido al de la versión secuencial como si lo era el de la versión centralizada. En esta estrategia, cada hilo cuenta con una seria fe \textit{buffers} que utiliza para la comunicación con el resto de hilos. Estos \textit{buffers} se implementan con las siguientes estructuras en los ficheros \verb|buffer.h| y \verb|buffer.c|:

\begin{lstlisting}
struct {
    int id;
    float gCost;
    int parent;
} buffer_elem_t;

struct {
    int capacity;
    int size;
    buffer_elem_t *elems;
} buffer_t;
\end{lstlisting}

Principalmente, cada hilo contará con un \textit{buffer} de recepción que el resto de hilos deberán boquear mediante un cerrojo para escribir en él los nodos que le pertenezcan, llamado \verb|incomebuffer|. Como hemos comentado en la Sección \ref{sec:descentralizado}, estos bloqueos sobre los \verb|incomebuffer| provocan una sobrecarga en la comunicación y disminuyen el rendimiento general del algoritmo. Para aliviar estos bloqueos, se ha seguido la siguiente estrategia: cada hilo cuenta con una serie de \textit{buffers} llamados \verb|outgobuffer|, uno por cada hilo, en los que irá incluyendo los nodos que le pertenezcan a otros hilos si en el momento de mandárselos no puede obtener el cerrojo correspondiente, ya sea porque otro hilo está añadiendo nodos, o porque el propietario lo está vaciando. Esta estrategia consigue que los hilos esperen a mandar los nodos cuando tienen el \verb|incombuffer| correspondiente disponible, pero puede provocar que un hilo haga trabajo innecesario, pues quien tiene que mandarle los nodos más prometedores no es capaz de obtener el cerrojo durante un largo periodo de tiempo. Para evitar esto, hemos establecido un número máximo de elementos, \verb|outgo_threshold|, que puede haber en cada \verb|outgobuffer|. Si el número de elementos alcanza este valor, el hilo no seguirá hasta no enviar los nodos a su propietario. Para conseguir este comportamiento, OpenMP nos proporciona el siguiente método:

\begin{lstlisting}
int omp_test_lock(omp_lock_t *lock);
\end{lstlisting}

\noindent que devuelve 1 (\textit{true}) si se logrado obtener el cerrojo y 0 (\textit{false}) en caso contrario. Este mecanismo también se aplica cuando un hilo comprueba si ha recibido nodos. Si alguno otro sigue escribiendo nodos en su \verb|incomebuffer|, éste no se bloquea y sigue su ejecución. A continuación podemos ver las primeras líneas del bucle principal (Código \ref{incomebuffer}), en el que se hace esta comprobación (líneas 2-14):

{\lstset{frame=trbl,
         numbers=left,
         caption=Comprobación del \textit{buffer} de recepción,
         label=incomebuffer}
\begin{lstlisting}
...
if (incomebuffers[tid]->size > 0) {
    omp_set_lock(&terminate_lock);
    terminate[tid] = 0;
    omp_unset_lock(&terminate_lock);
    if (incomebuffers[tid]->size >= income_threshold) {
        omp_set_lock(&incomebuffers_locks[tid]);
        fill_buffer(tmp_buffer, incomebuffers[tid]);
        omp_unset_lock(&incomebuffers_locks[tid]);
    } else if (omp_test_lock(&incomebuffers_locks[tid])) {
        fill_buffer(tmp_buffer, incomebuffers[tid]);
        omp_unset_lock(&incomebuffers_locks[tid]);
    }
}

if (tmp_buffer->size > 0) {
    for (int i = 0; i < tmp_buffer->size; i++) {
        buffer_elem_t msg = tmp_buffer->elems[i];
        if (visited[msg.node_id] != NULL) {
            if (msg.gCost < visited[msg.node_id]->gCost) {
                visited[msg.node_id]->gCost = msg.gCost;
                visited[msg.node_id]->fCost = msg.gCost + source->heuristic(msg.node_id, goal_id);
                visited[msg.node_id]->parent = msg.parent_id;
                if (visited[msg.node_id]->is_open) heap_update(open, visited[msg.node_id]);
                else heap_insert(open, visited[msg.node_id]);
            }
        } else {
            visited[msg.node_id] = node_create(msg.node_id, msg.gCost, msg.gCost + source->heuristic(msg.node_id, goal_id), msg.parent_id);
            heap_insert(open, visited[msg.node_id]);
        }
    }
    tmp_buffer->size = 0;
}
...
\end{lstlisting}
}

\noindent El valor de \verb|income_threshold| es el número máximo de elementos para forzar la recepción de nodos, al igual que en el envío. En este caso, para evitar bloquear el \textit{buffer} de recepción si otros hilos tienen que insertar nodos, se copia el contenido a un \textit{buffer} temporal, y posteriormente se tratan como si de la expansión en la versión secuencial se tratara (líneas 16-33). En este caso, al contrario que en la versión centralizada, las listas abierta y cerrada son privadas, por lo que esta segunda parte se hace sin necesidad de obtener ningún cerrojo.

A continuación, se realiza la comprobación de la condición de terminación de la misma forma que en la versión centralizada (\ref{terminacion}). Finalmente, se expande el nodo obtenido de la lista abierta local (Código \ref{expansion}). Para cada vecino se deberá calcular qué hilo es el propietario. Esta operación se implementa en HDA* mediante un función \textit{hash}:

\begin{lstlisting}
static inline int hash(int n_id, int k) {
    return n_id % k;
}
\end{lstlisting}

\noindent Si el nodo pertenece al propio hilo, se puede comprobar directamente si es duplicado o hay que insertarlo en las listas locales (líneas 6-18). En caso de pertenecer a otro hilo se sigue el esquema que se ha presentado antes: si el número de elementos a enviar es igual o mayor que el valor de umbral, se envían (líneas 19-23); si no lo supera pero ha podido obtener el cerrojo, también los envía (líneas 24-29); y si no es ninguno de los casos, los guarda en el \verb|outgobuffer| correspondiente para enviarlos después (líneas 30-32).

\clearpage

{\lstset{frame=trbl,
         numbers=left,
         caption=Evaluación de los nodos vecinos en HDA*,
         label=expansion}
\begin{lstlisting}
...
for (int i = 0; i < neighbors->count; i++) {
    int n_id = neighbors->nodeIds[i];
    float new_cost = current->gCost + neighbors->costs[i];
    int owner = hash(n_id, k);
    if (owner == tid) {
        if (visited[n_id]) {
            if (new_cost < visited[n_id]->gCost) {
                visited[n_id]->gCost = new_cost;
                visited[n_id]->fCost = new_cost + source->heuristic(n_id, goal_id);
                visited[n_id]->parent = current->id;
                if (visited[n_id]->is_open) heap_update(open, visited[n_id]);
                else heap_insert(open, visited[n_id]);
            }
        } else {
            visited[n_id] = node_create(n_id, new_cost, new_cost + source->heuristic(n_id, goal_id), current->id);
            heap_insert(open, visited[n_id]);
        }
    } else if (outgobuffers[owner]->size > outgo_threshold) {
        omp_set_lock(&incomebuffers_locks[owner]);
        buffer_insert(incomebuffers[owner], (buffer_elem_t){n_id, new_cost, current->id});
        fill_buffer(incomebuffers[owner], outgobuffers[owner]);
        omp_unset_lock(&incomebuffers_locks[owner]);
    } else if (omp_test_lock(&incomebuffers_locks[owner])) {
        buffer_insert(incomebuffers[owner], (buffer_elem_t){n_id, new_cost, current->id});
        if (outgobuffers[owner]->size > 0) {
            fill_buffer(incomebuffers[owner], outgobuffers[owner]);
        }
        omp_unset_lock(&incomebuffers_locks[owner]);
    } else {
        buffer_insert(outgobuffers[owner], (buffer_elem_t){n_id, new_cost, current->id});
    }
}
\end{lstlisting}
}

\section{Resultados}\label{sec:resultados}

A continuación vamos a evaluar los tres algoritmos en un conjunto de escenarios  y compararemos los tiempos obtenidos. La colección de ejemplos se ha obtenido de la página web del profesor Nathan R. Sturtevant\footnote{\href{https://movingai.com/}{https://movingai.com/}}\cite{sturtevant2012benchmarks}. Los escenarios de prueba consisten en \textit{grids} de mapas de videojuegos conocidos, como \textit{Dragon Age}, \textit{Baldur's Gate II} o \textit{Starcraft}, además de otros ejemplos sintéticos como laberintos o mapas aleatorios (Figura \ref{fig:mapas}).

\begin{figure}[H]
    \centering
    \begin{subfigure}[b]{0.32\textwidth}
        \centering
        \includesvg[width=\linewidth]{images/lak505d.svg}
        \caption{\textit{Dragon Age}}
        \label{dao}
    \end{subfigure}
    \hfill
    \begin{subfigure}[b]{0.32\textwidth}
        \centering
        \includesvg[width=\linewidth]{images/Cauldron.svg}
        \caption{\textit{StarCraft}}
        \label{starcraft}
    \end{subfigure}
    \hfill
    \begin{subfigure}[b]{0.32\textwidth}
        \centering
        \includesvg[width=\linewidth]{images/maze512-16-0.svg}
        \caption{Laberinto}
    \end{subfigure}
    \caption{Ejemplos de los mapas de prueba.}
    \label{fig:mapas}
\end{figure}

Es precisamente en este tipo de escenarios representados por \textit{grids} en los que las versiones paralelas de los algoritmos de búsquedas suelen tener más problemas en obtener una mejora con respecto a la versión secuencial. En el artículo de Burns et al. \cite{burns2010best} en el que comparan distintas versiones paralelas de la búsqueda primero el mejor, entre las que se incluyen PRA* y HDA*, obtienen una probabilidad de que un hilo no pueda acceder a un nodo, porque otro lo está modificando en ese momento, de un \% cuando se permiten movimientos en las cuatro direcciones cardinales, y de un 45\% cuando los movimientos se permiten en las ocho direcciones. Eso sumado a que en este tipo de problemas la operación de expansión es poco costosa, hace que los hilos pasen la mayor parte del tiempo en espera para obtener los cerrojos. Además, cuanto más cercano es el valor de la heurística al valor real, más tenderán los hilos a buscar todos por la misma zona.

Todos los tiempos que se mostrarán a continuación se han obtenido ejecutando los algoritmos sobre el mapa correspondiente en diez casos con un coste de camino similar y tomando la media de esas ejecuciones.

Vamos a empezar con un conjunto de mapas de tipo laberinto. En la Figura \ref{mazes} podemos ver los resultados para las ejecuciones según crece el número de hilos. Los dos primeros números del título en cada gráfica representan la altura y la anchura total del mapa, y el tercer número en este caso representa el ancho de los caminos del laberinto.

\begin{figure}[H]
    \centering
    \begin{subfigure}[b]{0.49\textwidth}
        \includesvg[width=\textwidth]{plots/maze/maze_1.svg}
    \end{subfigure}
    \hfill
    \begin{subfigure}[b]{0.49\textwidth}
        \includesvg[width=\textwidth]{plots/maze/maze_2.svg}
    \end{subfigure}

    \vspace{0.3em}
    
    \begin{subfigure}[b]{0.49\textwidth}
        \includesvg[width=\textwidth]{plots/maze/maze_4.svg}
    \end{subfigure}
    \hfill
    \begin{subfigure}[b]{0.49\textwidth}
        \includesvg[width=\textwidth]{plots/maze/maze_8.svg}
    \end{subfigure}

    \vspace{0.3em}

    \begin{subfigure}[b]{0.49\textwidth}
        \includesvg[width=\textwidth]{plots/maze/maze_16.svg}
    \end{subfigure}
    \hfill
    \begin{subfigure}[b]{0.49\textwidth}
        \includesvg[width=\textwidth]{plots/maze/maze_32.svg}
    \end{subfigure}

    \vspace{0.3em}

    \begin{subfigure}[b]{0.49\textwidth}
        \includesvg[width=\textwidth]{plots/synthetic/maze1000.svg}
    \end{subfigure}

    \caption{Resultados de tiempos de ejecución frente al número de hilos en mapas de tipo laberinto. El primer y segundo número en el título representan las dimensiones del mapa y el tercer número la anchura de los pasillos.}
    \label{mazes}
\end{figure}

De estos resultados podemos extraer varias conclusiones que veremos replicadas en el resto de ejemplos más adelante. En primer lugar, el tiempo de SPA* no deja de aumentar con el número de hilos en todos los casos. Esto se debe a la sobrecarga en la sincronización entre procesos, y no deja de ser otra prueba de que una paralelización directa del algoritmo empeora enormemente la eficiencia con respecto a la versión secuencial, lo que en la literatura se conoce como un acercamiento ingenuo (\textit{naive approach}). En cuanto a HDA* vemos que si que consigue alguna mejora con respecto a la versión secuencial en algunos casos. Estos casos dependen del número de hilos utilizados y de la estructura del mapa. En general se observa la tendencia a reducir el tiempo de ejecución cuando se aumentan el número de hilos, lo que lo convierte en un buen candidato a una formulación paralela y escalable del algoritmo. En algunos casos, como el primer laberinto, vemos que solamente consigue una mejora cuando se usan cinco hilos. Lo que ocurre es que HDA*, a costa de aliviar la sobrecarga de sincronización, introduce una sobrecarga de búsqueda y de comunicación. Esto puede ser determinante en el tiempo de ejecución, sobre todo cuando los mapas son pequeños. Como prueba de ello, podemos ver que cuando se utiliza un mapa más grande, como en el último caso, la mejora es mucho más notable, llegando a reducirse el tiempo de ejecución en un $35\%$ cuando se utilizan cuatro hilos.

En la Figura \ref{games} podemos apreciar cómo afecta el tamaño del mapa a la mejora obtenida con HDA*. Estos mapas son los mostrados en las Figuras \ref{dao} y \ref{starcraft} correspondientes a los videojuegos \textit{Dragon Age} y \textit{Starcraft}. Al lado del nombre entre paréntesis se muestran las dimensiones de cada mapa.

\begin{figure}[H]
    \centering
    \begin{subfigure}[b]{0.49\textwidth}
        \includesvg[width=\textwidth]{plots/games/dao.svg}
    \end{subfigure}
    \hfill
    \begin{subfigure}[b]{0.49\textwidth}
        \includesvg[width=\textwidth]{plots/games/starcraft.svg}
    \end{subfigure}

    \caption{Tiempos de ejecución frente al número de hilos involucrados en la búsqueda.}
    \label{games}
\end{figure}

Otro conjunto de mapas en el que se han probado estos algoritmos son mapas aleatorios (Figura \ref{random}). Consisten en cuadrículas de $512 \times 512$ en las que se colocan obstáculos de manera aleatoria. Contamos con siete escenarios correspondientes a diferentes porcentajes de ocupación por parte de los obstáculos, desde 10\% hasta un 40\%.

\begin{figure}[H]
    \centering
    \begin{subfigure}[b]{0.49\textwidth}
        \includesvg[width=\textwidth]{plots/random/random_10.svg}
    \end{subfigure}
    \hfill
    \begin{subfigure}[b]{0.49\textwidth}
        \includesvg[width=\textwidth]{plots/random/random_15.svg}
    \end{subfigure}

    \vspace{0.3em}
    
    \begin{subfigure}[b]{0.49\textwidth}
        \includesvg[width=\textwidth]{plots/random/random_20.svg}
    \end{subfigure}
    \hfill
    \begin{subfigure}[b]{0.49\textwidth}
        \includesvg[width=\textwidth]{plots/random/random_25.svg}
    \end{subfigure}

    \vspace{0.3em}

    \begin{subfigure}[b]{0.49\textwidth}
        \includesvg[width=\textwidth]{plots/random/random_30.svg}
    \end{subfigure}
    \hfill
    \begin{subfigure}[b]{0.49\textwidth}
        \includesvg[width=\textwidth]{plots/random/random_35.svg}
    \end{subfigure}

    \vspace{0.3em}

    \begin{subfigure}[b]{0.49\textwidth}
        \includesvg[width=\textwidth]{plots/random/random_40.svg}
    \end{subfigure}

    \caption{Resultados de tiempos de ejecución frente al número de hilos en mapas de tipo aleatorio de dimensiones $512 \times 512$.}
    \label{random}
\end{figure}

En este tipo de mapas se observa una mejora mucho mas estable que en los mapas de laberintos cuando el número de hilos aumenta. En este caso el espacio de búsqueda se reparte mejor entre los distintos hilos, llegando incluso a reducir a la mitad el tiempo de búsqueda en algunos casos.

Finalmente, vamos a evaluar el rendimiento de los tres algoritmos en unos mapas sintéticos para mostrar cómo afecta la naturaleza y la estructura del espacio de búsqueda. Se han construido dos escenarios, uno de tamaño $5000 \times 5000$ sin ningún obstáculo, en el que el nodo de inicio se encuentra arriba a la izquierda y el objetivo abajo a la derecha; y otro de tamaño $954 \times 7$ en el que solamente hay dos caminos hacia el objetivo, ambos con el mismo coste. Podemos ver una representación de éste último en la Figura \ref{fig:custom}, donde el círuclo verde representa el inicio y el rojo el objetivo.

\begin{figure}[H]
    \centering
    \includesvg[width=0.5\linewidth]{plots/synthetic/test.svg}
    \caption{Escenario de prueba con dos caminos óptimos para alcanzar la solución.}
    \label{fig:custom}
\end{figure}

En la Figura \ref{synthetic} podemos ver como ambos algoritmos funcionan peor en estos casos que la versión secuencial, además se puede observar que HDA* funciona peor que SPA* para estos escenarios. 

\begin{figure}[H]
    \centering
    \begin{subfigure}[b]{0.49\textwidth}
        \includesvg[width=\textwidth]{plots/synthetic/test.svg}
    \end{subfigure}
    \hfill
    \begin{subfigure}[b]{0.49\textwidth}
        \includesvg[width=\textwidth]{plots/synthetic/noobstacle.svg}
    \end{subfigure}

    \caption{Resultados de tiempos de ejecución frente al número de hilos en mapas sintéticos. El primero corresponde al escenario de la Figura \ref{fig:custom} y el segundo a un escenario sin ningún obstáculo.}
    \label{synthetic}
\end{figure}

El mal desempeño de HDA* en el mapa de dos caminos se debe a que los hilos pasan la mayor parte del tiempo esperando a recibir un nodo para expandir. Por cómo se ha implementado la función \textit{hash} para repartir los nodos entre hilos, cada nodo consecutivo del camino óptimo pertenecerá a un proceso distinto, por lo que solo el hilo con el nodo actual en el camino óptimo estará realmente trabajando. Cabe destacar que podría obtenerse un mejor resultado diseñando una función \textit{hash} específica para este escenario que repartiera el trabajo entre los hilos de otra manera, como por ejemplo asignando cada fila a un hilo. En el caso del mapa vacío, lo que ocurre es que la heurística guía al camino perfectamente, por lo que la versión secuencial solamente explora los nodos en el camino óptimo, y cualquier otro trabajo extra realizado por HDA* afecta negativamente a su resultado, además del coste añadido por comunicar los nodos entre procesos.

A continuación, vamos ver cómo afecta al rendimiento de HDA* el parámetro \verb|threshold|, que recordemos que es el número máximo de nodos que puede acumular un hilo antes de enviárselos a su propietario. En realidad definíamos dos valores: uno para el \textit{buffer} de recepción y otro para el de envío. Asignamos el mismo valor a ambos, por eso en la Figura \ref{fig:threshold} se nombra el eje de abscisas como \verb|threshold|. Se ha ejecutado el algoritmo con cuatro y ocho hilos en el mapa con forma de laberinto de $1000 \times 1000 \times 1$.

\begin{figure}[H]
    \centering
    \begin{subfigure}[b]{0.49\textwidth}
        \includesvg[width=\textwidth]{plots/threshold_4.svg}
    \end{subfigure}
    \hfill
    \begin{subfigure}[b]{0.49\textwidth}
        \includesvg[width=\textwidth]{plots/threshold_8.svg}
    \end{subfigure}

    \caption{Efecto del parámetro \texttt{threshold} en el algoritmo HDA* con cuatro y ocho hilos. Eje de abscisas en escala logarítmica.}
    \label{fig:threshold}
\end{figure}

Se observa una tendencia a disminuir el tiempo de ejecución a medida que aumenta el valor de este parámetro, pero tambien se ve que aumentarlo demasiado puede ser contraproducente, debido a que algún hilo deba esperar demasiado a recibir los nodos que le llevarán a la solución óptima. Este comportamiento nos sugiere que existe un valor óptimo para este parámetro, y que probablemente dependerá de las características del problema. En la Sección \ref{sec:futuro} propondremos la búsqueda automatizada de este valor óptimo como una posible línea de investigación futura.

\begin{figure}[h]
\centering
\includesvg[width=0.49\textwidth]{plots/slowexp.svg}
\caption{Evolución del tiempo de ejecución en un problema con una operación de expansión costosa.}
\label{fig:slowexp}
\end{figure}

Por último, vamos a simular un problema en el que la operación de expansión es costosa. Esto se ha conseguido añadiendo una instrucción \verb|sleep(1)| en cada llamada a la función \verb|get_neighbors|. En la Figura \ref{fig:slowexp} podemos observar como en este tipo de escenarios ambas estrategias consiguen mejorar enormemente a la versión secuecnial. De hecho incluso SPA* pasa ser más prometedor que HDA*, obtieniendo un menor tiempo de ejecución y una mayor consistencia en la mejora a medida que aumentan los hilos.

\newpage
%%%%%%%%%%%%%%%%%%%%%%%%%%%%%%%%%%%%%%%%%%%%%%%%%%%%%%%%%%%%%%%%%%%%%%%%%%%%%%%
\section{Conclusiones} \label{sec:conclusiones}
%%%%%%%%%%%%%%%%%%%%%%%%%%%%%%%%%%%%%%%%%%%%%%%%%%%%%%%%%%%%%%%%%%%%%%%%%%%%%%%

A lo largo de este trabajo hemos estudiado las distintas estrategias de paralelización del algoritmo A* propuestas en las últimas décadas, destacando los enfoques centralizados y descentralizados. Se han analizado las ventajas e inconvenientes que presentan cada uno de ellos, así como los retos que surgen a la hora de implementarlos. Con el objetivo de evaluar ambos enfoques se han desarrolado tres implementaciones del algoritmo A*: una versión secuencial que ha servido como base para la comparación, una versión centralizada conocida como SPA* y una versión descentralizada adaptada a un entorno de memoria compartida basada en HDA*, uno de los resolvedores del estado del arte actual.

Tras estudiar el comportamiento de estas dos estrategias de paralelización, hemos visto rápidamente que la centralizada resulta ser muy ineficiente, empeorando a la versión secuencial a medida que se involucran un mayor número de hilos en la búsqueda. Por otro lado, con la versión distribuida, sí que se ha obtenido una mejora consistente, observando una tendencia a disminuir el tiempo de ejecución a medida que aumenta el número de hilos, llegando a duplicar la velocidad del algoritmo secuenial en algunos escenarios. Esta mejora ha sido más notable cuanto mayor era la complejidad del problema en términos de tamaño y ramificación del grafo. Sin embargo, en problemas pequeños hemos comprobado que esta versión se comporta mucho peor que la secuencial, incluso que la centralizada, debido a la sobrecarga introducida tanto en la búsqueda, por expandir nodos innecesarios, como de comunicación, por tener que repartir el trabajo entre hilos.

De la implementación que se ha propuesto en este trabajo para la versión descentralizada, se ha estudiado el efecto del parámetro \verb|threshold| que controla el margen con el que los hilos pueden seguir expandiendo nodos antes de enviarlos a su propietario. Se ha visto que un valor adecuado para este parámetro puede mejorar enormemente el resultado obtenido con esta estrategia. Finalmente, se ha probado a aplicar los algoritmos en un entorno en el que la operación de expansión resulta costosa, observando que en este tipo de escenarios las estrategias centralizadas siguien siendo una alternativa que vale la pena considerar.

Una conclusión general que podemos obtener más allá de los resultados específicos para cada enfoque es que el comportamiento del algoritmo A* y, en general, de sus versiones, depende enormemente de la estructura del grafo de búsqueda, lo cual explica por qué han surgido a lo largo de los años tantas versiones diferentes de un mismo algoritmo, cada uno con un ámbito de aplicación en el que destaca frente a otras. La paralelización de este algoritmo sigue siendo un área de ingestigación activa a día de hoy para intentar conseguir implementaciones cada vez más eficientes y escalables.

\newpage

%%%%%%%%%%%%%%%%%%%%%%%%%%%%%%%%%%%%%%%%%%%%%%%%%%%%%%%%%%%%%%%%%%%%%%%%%%%%%%%
\section{Trabajo futuro} \label{sec:futuro}
%%%%%%%%%%%%%%%%%%%%%%%%%%%%%%%%%%%%%%%%%%%%%%%%%%%%%%%%%%%%%%%%%%%%%%%%%%%%%%%

El trabajo desarrollado se ha centrado en la evaluación de las distintas estrategias de paralelización del algoritmo en un entorno de memoria compartida. Las conclusiones obtenidas nos han permitido identificar bajo qué condiciones y en qué tipo de escenarios resultan ser más prometedores cada una estas técnicas. A continuación proponemos una serie de ideas para motivar futuros trabajos en esta línea:

\begin{itemize}
    \item Analizar la eficiencia de las impementaciones propuestas en términos de uso de memoria y expansión de nodos para evaluar cuantitativamente la sobrecarga introducida en la búsqueda.
    \item Investigar más a fondo el efecto del parámetro \verb|threshold| en el rendimiento de HDA*, concretamente, automatizando su búsqueda y estudiando cómo afecta a la mejora obtenida con respecto a la versión secuencial en función de la estructura del grafo de búsqueda.
    \item Explorar las estrategias descentralizadas en entornos de memoria distribuida implementadas mediante paso de mensajes.
    \item Estudiar las nuevas implementaciones propuestas para arquitecturas GPU, en las que se aprovecha el paralelismo inherente a estos dispositivos para aumentar el número de expansiones.
\end{itemize}

Como se ha expuesto en la introducción, el tema elegido para este trabajo ha sido motivado por la realización de un Máster en Computación Cuántica en la Universidad Internacional de La Rioja, en el que se propuso una línea para un Trabajo de Fin de Máster enfocada en una propuesta cuántica del algoritmo A*. La motivación inicial fue investigar a fondo el desarrollo del algoritmo en este trabajo, y explorar las implementaciones paralelas propuestas en la literatura para, en un futuro, poder realizar una comparación con una versión cuántica del mismo. Hasta el momento, en relación a la propuesta cuántica, hemos revisado artículos como el de Montanaro \cite{montanaro2015quantum}, en el que propone una mejora del algoritmo de backtracking utilizando lo que se conoce como caminatas cuánticas, el equivalente cuántico a las caminatas aleatorias. También hay otros estudios más recientes como el de Li \cite{li2025exponential} en el que se trata específicamente el problema del camino más corto y se propone un algoritmo cuántico que consigue, de manera teórica, una ventaja exponencial sobre cualquier algoritmo clásico para resolver este tipo de problema en un tipo de grafos conocidos como grafos soldados. A día de hoy no hay ninguna propuesta cuántica del algoritmo A* como tal, y los avances en este tipo de problemas en computación cuántica son muy recientes.

\newpage

%%%%%%%%%%%%%%%%%%%%%%%%%%%%%%%%%%%%%%%%%%%%%%%%%%%%%%%%%%%%%%%%%%%%%%%%%%%%%%%%

\bibliography{bibfile}


%%%%%%%%%%%%%%%%%%%%%%%%%%%%%%%%%%%%%%%%%%%%%%%%%%%%%%%%%%%%%%%%%%%%%%%%%%%%%%%%

\appendix


\end{document}
